\section{C. Smooth monodromy and the seven-dimensional invariant}

The actual positive-loop smooth monodromy is zero. C.3.24 proves
this through an equivariant orbit-knot integer; C.3.25--C.3.27
give a second determination through an explicit spin$^c$
eight-dimensional coboundary and direct index additivity. Their
boundary-invariant and one-loop arguments are independent; their
actual annular representative and characteristic inputs are shared
and identified below. B.4.33 excludes the specific analytic
filling of Part B, so C.5.1 records that its requested Chern
data are undefined. In this section $M_\Phi$ is a closed seven-manifold,
whereas $W_8$ denotes a
compact oriented eight-manifold with boundary. The circle parameter
has positive orientation. On $M_\Phi$ we use the orientation
$d\theta\wedge o_{X_c}$; this is the boundary orientation of a
holomorphic family over a positively oriented disc. The statements
below do not assign a characteristic number to any unconstructed
filling in Part B.

\subsection{C.1. Clutching conventions and the mapping torus}

\begin{proposition}[C.1.1 (Known, with derivation): the mapping-torus comparison]
For an orientation-preserving smooth diffeomorphism
$\phi:S^6\to S^6$, put
\[
 M_\phi=(S^6\times[0,1])/((x,1)\sim(\phi(x),0)),\qquad
 \Sigma_\phi=D^7_+\cup_\phi(-D^7_-).
\]
The boundary identification uses the usual oriented parametrizations
of the two underlying copies of $S^6$; the minus sign reverses the
orientation of the second disc. Then
\[
 M_\phi\cong (S^1\times S^6)\mathbin\#\Sigma_\phi
\]
as oriented smooth manifolds. With the corresponding sphere-gluing
identification, $[\phi]\mapsto[\Sigma_\phi]$ identifies
$\pi_0\Diff^+(S^6)$ with $\Theta_7$. For either spin structure on
$M_\phi$,
\[
 \mu(M_\phi)=\mu(\Sigma_\phi).
\]
For the family in the prompt the gluing map is
$\Phi=F_{2,c_0-2}\circ H$, with $c(r)=c_0-2r$.
\end{proposition}

\begin{proof}
An ambient isotopy first arranges that $\phi$ is the identity on a
small closed disc: move the image of its centre back, correct the
oriented derivative through $\GL^+(6,\R)$, and use isotopy extension
on a smaller coordinate disc. Consequently its support lies in the
complementary disc $D^6$. Cut the mapping torus at its seam and take
the block $D^6\times[-\epsilon,\epsilon]$. The lower half of this
block has the old coordinates; on the upper half replace the
coordinates by $\phi^{-1}$. This identifies the reglued block with a
standard seven-disc. Its attaching map is $\phi$ on its upper
hemisphere and the identity on the lower hemisphere and side collar.
The upper hemisphere has boundary orientation $o_{D^6}$ for the
orientation $dt\wedge o_{D^6}$ on the block. Thus the replacement
has exactly the sign in $D^7_+\cup_\phi(-D^7_-)$.
Removing a small ball from the second disc in that sphere leaves the
same attaching construction; this proves the connected-sum formula.

Sphere gluing identifies pseudo-isotopy classes with homotopy spheres;
Cerf's pseudo-isotopy theorem changes pseudo-isotopy to isotopy for
the simply connected six-sphere. The exact primary statements are
Theorem 0 and Corollaries 1--2 of \citet[p.~8]{Cerf1970};
the sphere case is also explained in Theorem 1 and Corollary 1 of
\citet[p.~156]{Poenaru1971}. These statements compare the
sphere gluing and extension-over-a-disc constructions. This is the
high-dimensional theorem, not the separate statement $\Gamma_4=0$.
The oriented connected-sum group is Theorem 1.1 and Lemma 2.1 of
\citet[pp.~504--505]{KervaireMilnor1963}; their dimension table on
p.~504 gives order 28 in dimension seven.

The sphere summand is spin. The product $S^1\times S^6$ is spin and
has two spin structures; both bound $S^1\times D^7$ with the
corresponding choice on $S^1$. That coboundary has zero fourth
cohomology and zero signature, hence zero Eells--Kuiper invariant.
Boundary connected sum of this coboundary with a spin coboundary of
$\Sigma_\phi$ proves additivity directly: the fourth intersection
forms are an orthogonal sum and the Pontryagin evaluations add.
Thus the two boundary spin structures give the same value. This is
also the connected-sum property in \citet[Section 3]{EellsKuiper1962}.

Finally $e^{-\pi i(c_0-2r)}=e^{-\pi ic_0}e^{2\pi ir}$.
Parallel transport therefore reaches $X_{c_0-2}$ along the positive
$w$-circle; its endpoint is identified by $F_{2,c_0-2}$.
This confirms the sign and endpoint in R12. Two smooth connections
can be joined by a smooth path of horizontal distributions, producing
an isotopy of the resulting monodromies. Holomorphic automorphisms in
the connected group $\C^*$ also change the result only by isotopy.
\end{proof}

\begin{proposition}[C.1.2 (Known, with derivation): the positive generator]
The generator $1\in\Z/28$ is represented by the oriented sphere
bundle of a rank-four oriented real vector bundle on $S^4$ with
$e=1$ and $p_1=6$, oriented as the boundary of its disc bundle.
Its Eells--Kuiper invariant is $1/28$.
\end{proposition}
\begin{proof}
The clutching parameters $(h,j)=(2,-1)$ in the oriented
$S^3$-bundle construction give $e=h+j=1$ and $p_1=2(h-j)=6$;
we orient the base and fibre to give these signs. These are the
bundle conventions and computations of \citet[Section 6,
equations (12)--(13)]{EellsKuiper1962}. The Gysin sequence gives
the integral homology of $S^7$, and the homotopy sequence gives
simple connectivity. Thus the total space is a homotopy sphere.
Its disc bundle has intersection form $(1)$ on the zero section,
is spin because its second cohomology with $\Z/2$ coefficients
vanishes, and has $p_1^2=36$ and signature 1. Hence
\[
 \mu=\frac{36-4}{896}=\frac1{28}.
\]
The order-28 classification identifies this with the requested
positive generator. The integer and rational arithmetic is checked
in \path{code/c_exact.py}.
\end{proof}

\subsection{C.2. The corrected hypotheses for the characteristic formula}

\begin{theorem}[C.2.1 (Known, with derivation): Lead 6's formula]
Let $M=M_\phi$ have the orientation above. Let $W_8$ be a smooth
compact oriented spin$^c$ eight-manifold with $\partial W_8=M$.
Let $c\in H^2(W_8;\Z)$ be the first Chern class of its determinant
line, and choose a trivialization on the boundary which induces the
spin structure used on $M$. Write $p=p_1(TW_8)$, choose a relative
lift $\bar c$, and let $\bar p$ be the unique rational relative lift
of $p$. Then, with all products evaluated on $[W_8,M]$,
\begin{equation}\label{eq:c-spin-c}
 \mu(M)\equiv
 \frac{\bar p^{,2}-4\sigma(W_8)}{896}
 -\frac{\bar p\bar c^{,2}}{192}
 +\frac{\bar c^{,4}}{384}\pmod{\Z}.
\end{equation}
The value is independent of the lift $\bar c$. If $W_8$ is almost
complex, use its canonical spin$^c$ structure, for which
$c=c_1(TW_8)$ and $p=c_1^2-2c_2$. If $W_8$ is spin and its
determinant line is trivial, this reduces to
\[
 28\mu(M)=\frac{p_1^2-4\sigma(W_8)}{32}\pmod{28}.
\]
All the boundary cohomology and spin-structure conditions in this
statement hold for the mapping torus; the existence of the particular
geometric $W_8$ proposed in Part B is a separate assertion.
\end{theorem}

\begin{proof}
The mapping-torus exact sequence, or Proposition C.1.1, gives
$H^k(M;\Z)=0$ for $2\leq k\leq5$ and $H^1(M;\Z)=\Z$.
The same vanishing in degree two holds with $\Z/2$ coefficients,
so $M$ is spin. Its spin structures form a torsor for
$H^1(M;\Z/2)=\Z/2$. The exact sequence of the pair gives an
isomorphism
$H^4(W_8,M;\Q)\to H^4(W_8;\Q)$ and a surjection
$H^2(W_8,M;\Z)\to H^2(W_8;\Z)$, proving the assertions
about lifts and determinant trivializations.

If two lifts of $c$ differ by $\delta\alpha$ with
$\alpha\in H^1(M;\Z)$, the relative cup-product identity gives
$\delta\alpha\smile b=\delta(\alpha\smile i^*b)$.
The restrictions of $c$ and $p$ vanish on $M$. Every term in the
change of $\bar p\bar c^2$ or $\bar c^4$ is therefore zero;
in particular $(\delta\alpha)^2=0$. This proves lift independence.
There is no assertion here that the relative class $\bar c$ itself
is unique.

For a closed spin$^c$ eight-manifold the degree-eight part of the
Dirac index density is
\[
 [e^{c/2}\widehat A]_8
 =\frac{7p_1^2-4p_2}{5760}
       -\frac{p_1c^2}{192}+\frac{c^4}{384}.
\]
Eliminate $p_2$ using $\sigma=(7p_2-p_1^2)/45$.
The result is exactly the right side of
\eqref{eq:c-spin-c}. To justify its use on a boundary, glue $W_8$
to the negative of a spin coboundary inducing the same boundary spin
structure; extend the determinant line by its chosen trivialization.
The glued manifold is spin$^c$, so its Dirac index is an integer.
Novikov additivity and the relative-product comparison identify that
integer with the difference between \eqref{eq:c-spin-c} and the
spin formula for $M$. The spin formula is
\citet[Section 6, equation (11)]{EellsKuiper1962}.
The spin$^c$ version is also explicitly recorded, with normalization
$28\mu$ and characteristic class
$\widehat p=(p_1+c^2)/2$, by
\citet[Section 2.2, equation (2.1)]{CrowleyNordstrom2020}.
Substituting this $\widehat p$ gives the same coefficients; both
polynomial identities are verified exactly in \path{code/c_exact.py}.

A spin coboundary can be supplied using the connected-sum description:
bound the product by $S^1\times D^7$ and bound the sphere summand
by its spin coboundary. Thus no inference from simple connectivity of
$M$ is made; indeed $\pi_1(M)=\Z$. The original Eells--Kuiper
paper's unstructured formulation has an extra condition on
$H^1(W;\Z/2)\to H^1(M;\Z/2)$ to ensure spin structures glue.
Our explicit requirement that the boundary spin structures match
supplies precisely that gluing datum. By Proposition C.1.1 the two
possible spin structures on this particular $M$ have equal invariant.
Finally the splitting principle gives
$p_1(TW_8)=c_1^2-2c_2$ for an almost complex tangent bundle.
\end{proof}

\subsection{C.2 supplement. A computed coboundary with a different boundary}

\begin{proposition}[C.2.2 (Proved): the untwisted intersection form]
Let $W_A$ be a sufficiently small compact tubular representative of the
smooth total space $\mathcal X^+$ of A.2.2 over a closed $u$-disc.
Its boundary is the mapping torus of Engel's \emph{untwisted} family
$X(u)$. It is not the six-sphere mapping torus $M_\Phi$.
Write $\mathcal L=H^2(S;\Z)$, $F=c_1\mathcal O_S(S_\infty)$,
$m=c_1(L)$, and $v=6P-O-3F$ in the conventions of A.4.1.
Let $H$ be the infinity-section class on $R$ and $A$ the exceptional
Cartier divisor pulled back to $R^+$. Then
\[
 H_4(W_A;\Z)=
 \frac{\mathcal L\oplus\Z H\oplus\Z A}{\Z(m-A)}.
\]
For $\alpha=(a,h,b)\in H^2(R^+;\Z)$ set
\[
 d(\alpha)=(I-T^*)a+h(m-F)+bF.
\]
The middle intersection form of $W_A$ is
\[
 B(\alpha,\beta)=-d(\alpha)\cdot d(\beta).
\]
Its rank is two, its signature is $+2$, and its two nonzero Smith
factors are $1,11$. On the image modulo its radical it is represented,
in the ordered basis $(v,m)$, by
\[
 \begin{pmatrix}6&1\\1&2\end{pmatrix}.
\]
\end{proposition}
\begin{proof}
The local semistable charts $xy=u$ give a tubular representative which
retracts to $Q^+$; its boundary is the nearby-fibre mapping torus. One
may equivalently choose a regular neighbourhood of $Q^+$ in the smooth
total space and round its boundary. The fibrewise plumbing identifies
that boundary with the inverse image of a small positive $u$-circle.
The normalization mapping cylinder and $H_3(S)=0$ identify its fourth
homology with the cokernel of
$H_4(S)\to H_4(R^+)$. In the divisor coordinates of A.4.2 the two
sections are $D'_0=H+m-A$ and $D_\infty=H$. Their difference is
$m-A$, proving the group presentation.

The normalization is an immersed complex hypersurface in $W_A$.
Its normal line has class $-D'_0-D_\infty$, because the whole central
divisor is the divisor of the base coordinate $u$. Put
$r_0\alpha=a+bF$ and $r_\infty\alpha=a-hm$; these are the
restrictions to the two sections. The self-intersection of two
four-cycles pushed forward from the normalization contains the normal
Euler term and the intersections between distinct local branches.
All transverse complex branch intersections have positive sign.
Consequently it is
\begin{align*}
 &-r_0\alpha\cdot r_0\beta-r_\infty\alpha\cdot r_\infty\beta\\
 &\quad+r_0\alpha\cdot T^*r_\infty\beta
       +T^*r_\infty\alpha\cdot r_0\beta\\
 &=-\bigl(r_0\alpha-T^*r_\infty\alpha\bigr)
       \cdot\bigl(r_0\beta-T^*r_\infty\beta\bigr).
\end{align*}
This is an intersection-of-cycles calculation in the actual smooth
fourfold, rather than an assignment of a form to an unconstructed
central fibre. The isometry property of $T$ and $T^*m=m-F$ give
the stated expression for $d$. Also $d(m-A)=0$, so it descends.

The translation formula in A.4.1 shows that the image of $d$ is
$\langle F,v,m\rangle$. The surjectivity of the three intersection
coordinates proved there shows equality integrally. Its radical is
$\Z F$, since $F^2=Fv=Fm=0$ and
$v^2=-6$, $m^2=-2$, $vm=-1$. The remaining form is the displayed
positive matrix, of determinant eleven and with an entry equal to
one. Integral row and column reduction gives Smith factors $1,11$.
These computations, including the full eleven-by-eleven form, are
independently checked in \path{code/e_push_C_attempts.py}.
\end{proof}

\begin{proposition}[C.2.3 (Proved): a relative Pontryagin number for $W_A$]
In the coordinates of C.2.2, $p=p_1(TW_A)$ evaluates as
\[
 p(a,h,b)=-4(aF)-28h.
\]
It vanishes on the radical of the middle form and therefore has a
rational relative lift. The square of any such lift is
\[
 \bar p^{\,2}=\frac{4512}{11}.
\]
These numbers do not compute $\mu(M_\Phi)$. In particular the
canonical determinant class of $W_A$ restricts nontrivially to its
own boundary, so C.2.1 cannot be applied with that determinant.
\end{proposition}
\begin{proof}
Let $h:R^+\to R$ be the modification, $C$ its nodal elliptic-fibre
centre, and $c=c_1(R)=F+2H+m$. The centre is a regular embedding of
codimension two, also at its node: its ideal there is $(X,zv)$.
Its normal bundle has degree zero, and $c\cdot C=0$. The blow-up
and small resolution give
\[
 c_1(R^+)=h^*c-A,\qquad
 A\,h^*D\,h^*E=0,\quad
 A^2h^*D=-C\cdot D,\quad A^3=0.
\]
These are the regular-embedding blow-up identities. The subsequent
small resolution is crepant, and $A$ is the pullback of a Cartier
divisor, so they remain valid on $R^+$. The normal line to the
immersed normalization has class
$h^*N_0+A$, where $N_0=-(2H+m)$.

For completeness, the required second Chern pairings do not assume
that $R^+$ is K\"ahler. Analytic Hirzebruch--Riemann--Roch applies to
a compact complex manifold; see \citet[Theorem~4.3, pp.~564--565]{AtiyahSinger1968IIIE}.
The modification has $Rh_*\mathcal O_{R^+}=\mathcal O_R$.
For the blow-up this follows locally from its hypersurface description
in the projective line bundle and the vanishing of the direct images
of $\mathcal O_{\mathbb P^1}(-1)$; for the small resolution it follows
from the formal neighbourhood of its $(-1,-1)$ rational curve.
Thus $\chi(h^*D)=\chi(D)$ for every line bundle $D$ on $R$.
The linear coefficient in the Riemann--Roch polynomial for $kD$
and the preceding intersection identities give
\[
 c_2(R^+)h^*D=c_2(R)D+C\cdot D.
\]
Moreover $\chi(\mathcal O(A))=\chi(\mathcal O)=1$: on the
exceptional projective line bundle over $C$, the restriction is its
relative $\mathcal O(-1)$, which has zero derived direct image.
Pulling through the small resolution does not change this Euler
characteristic. Applying Riemann--Roch to $A$ yields $c_2(R^+)A=0$,
because $A^3$, $c_1(R^+)A^2$, and $c_1(R^+)^2A$ all vanish.

The real tangent bundle of $W_A$ restricted to the normalization is
the sum of its real tangent bundle and this normal line. Therefore
\[
 h^*p_1(R)+h^*N_0^2-4[C]
\]
has the same pairings on pulled-back divisors as $p_1(TW_A)$, and
the pairing with $A$ is zero. Here $[C]$ denotes its degree-four
cohomology class on $R$. Since $H^2=-mH$, we have
$(2H+m)^2=m^2=-2[\mathrm{pt}]$. The surface $S$ has
$p_1(S)=-24[\mathrm{pt}]$. The ruled-bundle tangent sequence thus
gives
\[
 p_1(R)+N_0^2=p_1(S)+2(2H+m)^2=-28[\mathrm{pt}],
\]
and $C\cdot(a+hH)=aF$. This proves the covector formula.

Modulo $\Z F$ the coordinate of $d(a,h,b)$ along $v$ is $aF$
and the coordinate along $m$ is $h$. Thus the covector on the
nondegenerate quotient is $(-4,-28)$. It annihilates the radical,
which is precisely the condition for a rational relative lift.
Its square is
\[
 (-4,-28)\frac1{11}
 \begin{pmatrix}2&-1\\-1&6\end{pmatrix}
 \binom{-4}{-28}=\frac{4512}{11}.
\]
Different relative lifts differ by a class pairing trivially with
absolute classes, so this value does not depend on the choice.

Finally the canonical first Chern class on the nearby untwisted
fibre is $f^*c_1\mathcal O_B(1)$. It is nonzero: the surjectivity of
$a\mapsto(aF,av,am)$ in A.4.1 supplies an integral two-cycle with
evaluation one, and the Mayer--Vietoris relations in A.4.2 do not
alter that evaluation. Hence no boundary trivialization of this
determinant exists. Even the formal expression
$(\bar p^{\,2}-4\sigma(W_A))/896=79/176$ is only an arithmetic
combination for this different boundary, not the invariant requested
in Part C. A cobordism implementing the finite logarithmic transforms
and the cusp closing map would be required to transport the result.
\end{proof}

\subsection{C.3. What an equivariant twist does and does not determine}

\begin{proposition}[C.3.1 (Proved): the circle action and local shear]
The effective circle subgroup $\rho=\R/\Z\subset\C/\Z$ of R6
acts freely outside its fixed double curve. That curve is a sphere,
and its two complex normal weights are $1,-1$. A monodromy
representative may be chosen $\rho$-equivariant. In marked real
coordinates on a regular torus fibre, its linear part is
$I-2\widehat\delta\otimes\gamma$; its affine correction is the
translation supplied by R2 and depends on the trivialization.
\end{proposition}

\begin{proof}
On a regular torus the circle acts by translation in the primitive
lattice direction $\widehat\delta$, and is therefore free.
At a multiple fibre, apply the invariant covector $\gamma$ to the
equation that a circle translation agrees with the $j$-th affine
deck generator to the power $a$. Since the deck translation has
$\gamma$-component $\pm1/m_j$, equality forces
$a\equiv0\pmod{m_j}$. Primitivity of $\widehat\delta$ then
forces the circle element to be the identity. This checks freeness
on the quotient using the actual twist vectors, including the root
centre; the argument is not just a statement about the product cover.

In a cusp toric chart, the relevant primitive edge directions are
$e_1,e_2,e_1-e_2$, with $e_2=\widehat\delta$.
Their nonzero determinants with $e_2$ have absolute value one.
Thus stabilizers off the edge in direction $e_2$ are trivial. Along
that edge the circle is fixed; the two normal characters in the
adjacent unimodular cones are $1$ and $-1$ because the base monomial
is invariant. The opposite-edge identification gives the double
curve with its two distinct triple-point endpoints, namely a
$\PP^1$. This verifies the toric step of R6 used here.

The circle acts on the entire family and commutes with the endpoint
isomorphism $F_2$. Averaging a Riemannian metric over this compact
circle makes the orthogonal horizontal distribution invariant.
Its parallel transport and hence the monodromy commute with $\rho$.
On a local marked torus chart,
$\Pi_{c-2}=\Pi_c(I-2\widehat\delta\otimes\gamma)$.
Inserting the identity complex-linear part and translation of R2
gives the asserted affine formula. A locally chosen lift of its
translation must not be multiplied by a real interpolation parameter
without checking its period cocycle; that operation need not descend
to the torus bundle.
\end{proof}

\begin{proposition}[C.3.2 (Proved): a global scalar gauge would be trivial]
Let a smooth circle act on a simply connected smooth manifold $Z$.
If $g:Z\to S^1$ is smooth and invariant under the circle, then
$z\mapsto g(z)\cdot z$ is smoothly isotopic to the identity.
The primitive coordinate $\theta$ detecting $\gamma$ on a regular
fibre of $X_c$ has no continuous extension to $X_c$.
\end{proposition}

\begin{proof}
Since $Z$ is simply connected, $g$ lifts to a smooth real function
$a$ with $g=e^{2\pi ia}$. Invariance implies
$a(\lambda z)-a(z)\in\Z$. This difference is continuous in the
connected circle and is zero at its identity, so it is zero
everywhere. The maps $z\mapsto e^{2\pi ita(z)}\cdot z$ are
diffeomorphisms with inverses obtained by replacing $t$ by $-t$.
They give the isotopy. In the present application a loop representing
$\widehat\gamma$ in a regular torus fibre has degree one under
$\theta$. If $\theta$ extended to the simply connected $S^6$,
that degree would be zero. The contradiction shows why the local
twist in Lead 7 is not a global scalar-gauge calculation.
\end{proof}

\begin{proposition}[C.3.3 (Known, with derivation): equivariance alone is insufficient]
All 28 oriented mapping classes of $S^6$ admit representatives
equivariant for one fixed semifree circle action with fixed set $S^2$.
Consequently the circle action, its fixed set, and its normal weights
do not determine $[\Phi]$.
\end{proposition}

\begin{proof}
Use $S^6\subset\operatorname{Im}\mathbb H\oplus\mathbb H$.
The equivariant Blakers--Massey construction gives a generator
$\sigma$ of $\pi_0\Diff^+(S^6)$ equivariant for simultaneous
quaternionic conjugation. The smooth reentrance statement is
\citet[Theorem 2.1]{DuranRigas2007}; Section 3.2 identifies this
representative and its powers with all 28 classes. Restrict
conjugation to $q=e^{i\theta/2}$, regarding $q$ and $-q$ as the
same effective circle element. Its fixed vector space consists of
$p\in\R i$ and $w\in\R+\R i$, a real three-space, so the
fixed set on the sphere is $S^2$. Outside it, at least one of the
two $(j,k)$-planes is nonzero, and an effective rotation fixing it
is the identity. Thus the action is semifree. All powers of
$\sigma$ commute with it. Conjugating the complex coordinate on
one of the two real normal planes writes the characters as
$1,-1$; this representation calculation is checked exactly in
\path{code/c_exact.py}. These are statements about smooth circle
actions and normal representations, not a comparison of their
complex structures with $X_c$.
\end{proof}

\begin{proposition}[C.3.4 (Proved): information lost under the proposed base change]
Under $w=r^6$ the smooth monodromy class becomes $6[\Phi]$.
Its value alone leaves two possibilities for $[\Phi]$ in $\Z/28$,
differing by 14. In particular a vanishing value after that base
change would not prove that $[\Phi]=0$.
\end{proposition}
\begin{proof}
The induced homomorphism on fundamental groups sends a positive
$r$-circle to six positive $w$-circles, so its transport is
$\Phi^6$ up to isotopy. The kernel of multiplication by six on
$\Z/28$ is $\{0,14\}$. The explicit preimages of all elements
of its image are checked in \path{code/c_exact.py}.
\end{proof}

\begin{proposition}[C.3.5 (Known formula, with proved comparison): the closing translation]
Fix the standard twists and all root and cusp markings of R2.
There is a coherent choice of the even-shift maps
\(\widetilde F_{h,c}:X_c\to X_{c+h}\), \(h\in2\mathbb Z\),
with the following explicit ordinary-chart translation. Choose global
roots \(F^3=t\), \(G^4=t-1\) on the uniformizing cover and put
\(\ell=(2\pi i)^{-1}\log(F/G)\pmod{\mathbb Z}\), \(e=(0,1)^t\).
Then the translation is
\begin{equation}\label{eq:C-explicit-closing}
 b_{h,c}=-\frac h{12}\Pi_c\widehat\gamma
              +\left(h\ell-\frac{h^2}{24}\right)e.
\end{equation}
For \(h=12\), it is the scalar multiplication
\begin{equation}\label{eq:C-twelve-scalar}
 K(t)=\frac{t^4}{(t-1)^3}
\end{equation}
on the ordinary \(\mathbb C^*\)-coordinate, after removing an integral
period. On the finite root charts its scalar representatives are
\(K(t)s_1^{-12}\) and \(K(t)s_2^{12}\), respectively, both units.
In a cusp coordinate \(q=t_c\), the scalar representative is
\(q^{-1}H(q)\), with \(H\) a holomorphic nonvanishing function.
The same map has a holomorphic \((\mathbb C^*)^2\)-unit representative
on the entire cusp filling. These two representatives differ by
exactly the \(\widehat\gamma\) cusp deck transformation.

For six positive \(w\)-loops the closing map is
\(\widetilde F_{12,c_0-12}\); thus the stated cusp exponent is \(-1\)
with the orientation convention of C.1.1. This computes the closing
cocycle, not the mapping class of its composition with parallel
transport. In particular it does not imply \(6[\Phi]=0\).
\end{proposition}
\begin{proof}
The raw formula without \(-h^2e/24\), its three extension checks,
and the coherent correction occur in
\cite[\texttt{p2:prop:global-even-shift} and the subsequent
orbit-groupoid calculations]{DGKXv1push}. We verify their exact use
here; the source is version~1, and version~2 omits these statements.
The global roots exist because the zero orders of \(t\) and \(t-1\)
on the simply connected uniformizing surface are divisible by three
and four. The local characters and their coprime orders give
\[
 \ell(g_1z)-\ell(z)=-\frac13,\qquad
 \ell(g_2z)-\ell(z)=\frac14,\qquad
 \ell(g_0z)-\ell(z)=\frac1{12}\pmod{\mathbb Z}.
\]
A change of logarithm adds an integral \(h\widehat\delta\) period,
so \eqref{eq:C-explicit-closing} is a well-defined torus translation.
Since \(\Pi_{c+h}=\Pi_c(I+hE)\), the two ordinary lattices agree;
the source-to-target lattice map of the identity linear part is
\(I-hE\).

For compatible logarithm lifts, covariance gives the affine defects
\[
 b_{h,c}(g_jz)-R_j(z)b_{h,c}(z)
       =\Pi_c(g_jz)\eta_j,\qquad
 \eta_1=\frac h2(\widehat u-\widehat w-\widehat\delta),\quad
 \eta_2=\frac h2(-\widehat w+\widehat\delta).
\]
Their \(\gamma\)-coordinates vanish, so the same right sides use
\(\Pi_{c+h}\). The vectors are integral for even \(h\). At the
cusp the defect is zero: \((A_0-I)\widehat\gamma=-\widehat\delta\)
cancels the increment of \(\ell\). These equations prove ordinary
and cusp descent with the stated orientation of the peripheral
loops.

For the finite comparison put \(w_j=(2\pi i)^{-1}\log s_j\).
The functions \(F/(s_1G)\) and \(Fs_2/G\) are invariant
holomorphic units on the respective full root discs. Choose
holomorphic logarithms \(\kappa_j\) of them. Then
\[
 \ell=\epsilon_jw_j+\kappa_j\pmod{\mathbb Z},\qquad
 \epsilon_1=1,\quad\epsilon_2=-1.
\]
The exact finite-to-ordinary comparison satisfies
\(\sigma_{j,c+h}-\sigma_{j,c}=h\epsilon_jw_je\).
Consequently the required finite-chart translation is
\[
 b^{(j)}_{h,c}
 =-\frac h{12}\Pi_c\widehat\gamma
       +\left(h\kappa_j-\frac{h^2}{24}\right)e.
\]
It is holomorphic on the full root disc and satisfies
\[
 b^{(j)}_{h,c}(g_jz)-R_j(z)b^{(j)}_{h,c}(z)
       -\frac{h\epsilon_j}{m_j}e
 =\Pi_c(g_jz)\eta_j.
\]
This is precisely equivariance for the source and target affine
finite actions, including the changed period column. The exact
identity
\(\sigma_{j,c+h}+b^{(j)}_{h,c}
=\sigma_{j,c}+b_{h,c}\) proves overlap agreement. Thus no real
interpolation of a multivalued period has been used.

At the cusp let \(q=e^{2\pi i s}\) and
\(C_{c,21}=\beta^{\rm part}+c+s\), which is holomorphic at
\(q=0\). The function
\[
 \widehat H=e^{2\pi i s/12}F/G
 \quad\hbox{satisfies}\quad
 \widehat H^{12}=q\frac{t^4}{(t-1)^3}.
\]
The right side is a unit because \(t\) has a simple pole at the
cusp. The character formula makes \(\widehat H\) single-valued;
it therefore extends as a nonvanishing holomorphic function. Write
\(\ell=-s/12+\kappa_0(q)\), with \(\kappa_0\) holomorphic.
Since
\(\Pi_c\widehat\gamma=(6\mu,-s+C_{c,21})^t\),
the logarithmic terms cancel in \eqref{eq:C-explicit-closing}:
\[
 \bar b_{h,c}(q)=
 \left(-\frac h2\mu(q),\,
 -\frac h{12}C_{c,21}(q)+h\kappa_0(q)-\frac{h^2}{24}\right)^t.
\]
Exponentiating gives a holomorphic unit in \((\mathbb C^*)^2\).
On every toric chart its action multiplies a character coordinate
by an integral character of this unit, so it and its inverse
extend across every double curve and both triple points. The
source and target cusp lattice multipliers differ by
\((1,e^{2\pi i h\lambda_1})=(1,1)\), for
\(\lambda\in\mathbb Z^2\). They are therefore identical. The
unit multiplication commutes with their toric shears and descends.
Exponentiation of the displayed formula gives exactly the ordinary
translation on the punctured cusp. We have now verified all
three comparisons and their inverses.

The uncorrected translations have composition defect
\(-hk e/12\). The constants in
\eqref{eq:C-explicit-closing} cancel it, since
\((h+k)^2-h^2-k^2=2hk\). Hence
\[
 \widetilde F_{k,c+h}\widetilde F_{h,c}
       =\widetilde F_{h+k,c},\qquad
 \widetilde F_{0,c}=\mathrm{id}.
\]
Any difference from the fixed R2 choice is a scalar automorphism,
which lies in its connected \(\mathbb C^*\). The choice therefore
has the same smooth monodromy class.

For \(h=12\), remove the actual period
\(-\Pi_c\widehat\gamma\) from \eqref{eq:C-explicit-closing}.
The remaining translation is \((12\ell-6)e\). Its exponential
is \((F/G)^{12}=K(t)\), proving
\eqref{eq:C-twelve-scalar}. The finite corrections divide it by
\(s_1^{12}\) and multiply it by \(s_2^{12}\). The original
uniformizers are retained: these are units because
\(t/s_1^3\) and \((t-1)/s_2^4\) are units. They need not be
replaced by one.

Put \(H(q)=qK(t(q))=e^{24\pi i\kappa_0(q)}\). The cusp unit
for \(h=12\) is exactly
\[
 U_{12,c}(q)=
 \left(e^{-12\pi i\mu(q)},\,
       e^{-2\pi i C_{c,21}(q)}H(q)\right).
\]
The \(\widehat\gamma\) deck transformation has multiplier
\[
 \left(e^{12\pi i\mu(q)},\,
       e^{2\pi i C_{c,21}(q)}q^{-1}\right),
\]
because \(B_0(1,0)^t=(0,-1)^t\). Multiplying these two
expressions gives \((1,q^{-1}H)=(1,K)\). This proves the
claimed comparison, including its sign, on the entire cusp
quotient. In particular the scalar representative's pole is not
an obstruction to the holomorphic map itself: its unit
representative already extends.

Finally \(w=e^{-\pi i c}\) makes six positive loops the path
\(c(r)=c_0-12r\); its endpoint is closed by the positive shift
\(h=12\). A quotient isotopy proving an order bound would still
have to patch the ordinary scalar description with the two-coordinate
cusp unit description. Scaling their additive logarithms separately
produces on their overlap a fractional
\(\widehat\gamma\)-period, which is not an integral period during
the interpolation. This patching has not been established. The
exact substitutions, coherent correction, finite defects, and
cusp comparison are checked in
\path{code/d_push_monodromy.py} and
\path{code/d_push_closing.py}; these are symbolic checks of the
identities proved above.
\end{proof}
\begin{proposition}[C.3.6 (Proved): independence of the equivariant lift]
Let \(X=X_c\) carry the effective semifree circle action of C.3.1,
with fixed surface \(F\cong S^2\). If two smooth equivariant
diffeomorphisms of \(X\) induce the same map of its orbit space,
they are smoothly isotopic. More generally, a smooth isotopy of
orbit maps which is supported in a compact subset of the free
orbit stratum lifts to an equivariant isotopy. Thus, whenever the
annular orbit twist has been constructed with a smooth equivariant
lift, its class in \(\Theta_7\) is independent of that lift.
This assertion does not construct an isotopy of the annular twist.
\end{proposition}
\begin{proof}
Their quotient is an equivariant orbit-preserving diffeomorphism
\(h\). On \(X\setminus F\) it has a unique expression
\(h(x)=g(x)\cdot x\), with \(g\) a smooth invariant circle-valued
function. We check its extension across \(F\) directly. Choose an
equivariant tubular slice and trivialize its normal bundle locally
over \(F\). Conjugating the negative-weight complex coordinate
identifies the normal action with scalar multiplication on
\(\mathbb C^2\). Since \(h\) preserves each orbit, it fixes the
base variable \(y\) and has normal part
\[
 (f_1(y,z),f_2(y,z))=g(y,z)(z_1,z_2),\qquad z\ne0.
\]
In particular \(z_2f_1=z_1f_2\), an identity of smooth functions
also at zero. In the homogeneous Taylor expansion in the four real
normal variables, each pair \((p_1,p_2)\) therefore satisfies
\(z_2p_1=z_1p_2\). In the complexified polynomial ring with
variables \(z_1,z_2,\bar z_1,\bar z_2\), this implies
\(p_1=z_1q\) and \(p_2=z_2q\): first set \(z_1=0\) and divide
by \(z_1\), then cancel \(z_1\). This coefficient argument works
with coefficients depending smoothly on \(y\). It produces a
common formal Taylor series for \(g\).

The formal series can be realized by a smooth complex function
\(G(y,z)\). Here is the parametric Borel construction needed for
this assertion. On a relatively compact base slice, write its
homogeneous terms as \(q_n(y,z)\), choose a smooth cutoff equal to
one near zero, and replace each term by
\(\chi(z/\epsilon_n)q_n(y,z)\). Choose
\(\epsilon_n\downarrow0\) so rapidly that the term and all
mixed base/normal derivatives of total order at most \(n/2\)
are bounded by \(2^{-n}\), away from the finitely many lower
degrees. This is possible because each such derivative retains a
positive power of \(\epsilon_n\). The resulting series converges
in every smooth seminorm, and its Taylor series is the prescribed
one: the cutoff is constant near zero and higher homogeneous
terms have zero lower jets. This proves the required local
parametric extension. The remainders
\(f_j-Gz_j\) are flat in the normal variables, locally uniformly
with all derivatives in \(y\). For \(z\ne0\),
\[
 g-G=\frac{\bar z_1(f_1-Gz_1)+\bar z_2(f_2-Gz_2)}{|z|^2}.
\]
A smooth function flat along \(z=0\), divided by \(|z|^2\),
extends smoothly and flatly by zero: the usual quotient derivative
formula and Taylor bounds of arbitrarily high order prove this for
every derivative. Thus \(g\) extends smoothly across the slice.
Its absolute value is one by continuity. Equivariance implies
invariance on the free part, hence on the slice by continuity.
The local extensions agree because they agree on the dense free
stratum. We obtain a global smooth invariant \(g:X\to S^1\).

Since \(X\) is simply connected, the covering
\(\mathbb R\to S^1\), \(a\mapsto e^{2\pi ia}\), gives a smooth
lift \(a\). The integer-valued difference
\(a(\lambda x)-a(x)\) is zero at \(\lambda=1\) and continuous
in the connected circle, so \(a\) is invariant. The maps
\(x\mapsto e^{2\pi isa(x)}\cdot x\), \(0\leq s\leq1\),
are diffeomorphisms with the explicit inverse obtained by changing
\(s\) to \(-s\). They give the required isotopy.

For the second assertion the free quotient is a principal circle
bundle. A connection lifts the compactly supported time-dependent
vector field of the orbit isotopy to an invariant vector field.
Its flow exists for the full finite time interval, is equivariant,
and extends by the identity near \(F\). The endpoint differs from
any other smooth equivariant lift of the same orbit map by the
orbit-preserving diffeomorphism just treated.

The general orbit-preserving statement is Haefliger and Salem
(1991, Theorem~3.1). We used the direct slice argument because the
non-equivariant Lemma~3.2 in that paper, as originally printed,
requires correction: complex conjugation on \(\mathbb C\)
preserves scalar-circle orbits but its quotient
\(\bar z/z\) has no continuous extension to zero. The theorem
with equivariance is unaffected; the corrected argument is also
explained in Karshon's research slides, which reproduce the original
statement. No conclusion about smooth maps of the topological orbit
space near \(F\) is assumed here.
\nocite{HaefligerSalem1991D,KarshonOrbitSlidesD}
\end{proof}
\begin{proposition}[C.3.7 (Proved): the central cusp orbit fibre]
Let \(H=W/\rho\), and let \(F\cong S^2\) be the fixed double
curve. Then
\[
 \pi_1(H)=\mathbb Z^2,\qquad
 H_i(H;\mathbb Z)=
 \begin{cases}
 \mathbb Z&i=0,3,\\
 \mathbb Z^2&i=1,2,\\
 0&i>3.
 \end{cases}
\]
The class of \(F\) is nonzero and primitive in \(H_2(H;\mathbb Z)\).
In particular \(F\) does not bound a three-chain in the central
orbit fibre. This statement does not assert a deformation retraction
of the entire cusp neighbourhood onto \(H\), or compute the class
of the annular orbit twist.
\end{proposition}
\begin{proof}
Number the six toric boundary curves of the normalization
\(\widetilde W=\mathrm{dP}_6\) cyclically as
\(D_1,\ldots,D_6\), with \(D_3,D_6\) the fixed pair, and put
\(p_i=D_i\cap D_{i+1}\), with indices modulo six. Opposite sides
are identified by the maps already checked in R6 and C.3.1.
Their endpoint rule is
\[
 p_{i-1}\longmapsto p_{i+3},\qquad
 p_i\longmapsto p_{i+2}\quad(D_i\longrightarrow D_{i+3}).
\]
It follows either from the two triangle classes of the periodic
fan or directly from these six endpoint equations that the vertex
classes are
\(A=\{p_2,p_4,p_6\}\) and \(B=\{p_1,p_3,p_5\}\).
The multiplicative units in the side maps preserve this endpoint
rule. The fixed curve is their resulting sphere, with endpoints
\(A,B\).

Put \(Y=\widetilde W/\rho\). This is an oriented compact
three-manifold whose boundary consists of the two fixed spheres.
Here are the local checks, including the two remaining vertices.
At a free orbit the quotient is an oriented three-dimensional
slice. At a fixed divisor the local model is
\(\mathbb C\times\mathbb C\), with the circle rotating only
the first factor, and its quotient is a half-space. The two toric
vertices not on \(D_3\cup D_6\) are isolated fixed points with
weights \((1,-1)\). Conjugating one coordinate gives the scalar
Hopf action on \(\mathbb C^2\), whose quotient is the cone on
\(\mathbb CP^1=S^2\), namely \(\mathbb R^3\). Thus they are
interior points, not additional boundary components.
All other stabilizers are trivial: the determinants of the circle
cocharacter with the four nonfixed edge rays are \(\pm1\).

The quotient is simply connected. Indeed local slices give path
lifts of orbit-space paths, and lifts on successive subintervals
can be joined by paths in their connected circle orbits. A lifted
loop closes after adding a path in its final orbit; its projection
is the original loop with a constant segment. Since
\(\widetilde W\) is simply connected, the resulting loop is
null-homotopic, and projecting the null-homotopy proves the claim.
Poincar\'e--Lefschetz duality and the boundary sequence now give
\[
 H_1(Y)=H_3(Y)=0,\qquad H_2(Y)=\mathbb Z.
\]
More explicitly, \(H_3(Y,\partial Y)=\mathbb Z\) maps to
\(H_2(\partial Y)=\mathbb Z^2\) by the sum of the two boundary
fundamental classes, and \(H_2(Y,\partial Y)=H^1(Y)=0\).
No classification theorem for simply connected three-manifolds
is needed for these conclusions.

Let \(\widetilde D\subset Y\) be the quotient of the six
boundary curves, and let \(D\subset H\) be its image. The space
\(\widetilde D\) is two spheres joined by two paths, each
subdivided into two edges. Hence
\[
 H_0(\widetilde D)=\mathbb Z,\quad
 H_1(\widetilde D)=\mathbb Z,\quad
 H_2(\widetilde D)=\mathbb Z^2.
\]
The image \(D\) is the sphere \(F\) together with two arcs
\(E_1,E_2\), each joining \(A\) to \(B\). Consequently
\[
 H_0(D)=\mathbb Z,\quad H_1(D)=\mathbb Z^2,\quad
 H_2(D)=\mathbb Z.
\]
Choose a path \(c\subset F\) from \(A\) to \(B\). The
endpoint rule gives the image of the generator of
\(\pi_1(\widetilde D)\) as
\[
 E_1E_2^{-1}cE_1^{-1}E_2c^{-1}.
\]
Writing \(a=E_1c^{-1}\), \(b=E_2c^{-1}\), this is
\(ab^{-1}a^{-1}b\). In particular the induced map on first
homology is zero.

We also fix the second-homology sign. Near a fixed divisor use a
complex normal coordinate \(re^{i\theta}\) and a complex tangent
coordinate. Orient the quotient by putting the positive circle
direction first in the orientation of the surface. At circle
weight \(+1\) the quotient boundary orientation is the complex
curve orientation; at weight \(-1\) it is its negative. The
two normal weights are opposite. Thus, although their boundary
orientations sum to zero in \(H_2(Y)\), the two fixed curves
with their complex orientations represent the same generator
of \(H_2(Y)\). Their maps to \(F\) are holomorphic of degree
one. Therefore the Mayer--Vietoris map
\[
 H_2(\widetilde D)\longrightarrow H_2(Y)\oplus H_2(D)
\]
has matrix
\[
 \begin{pmatrix}1&1\\-1&-1\end{pmatrix}.
\]
It has a primitive rank-one image, rank-one kernel and torsion-free
rank-one cokernel.

The normalization conductor is a topological pushout, and taking
a compact group quotient commutes with that pushout. The inclusion
\(\widetilde D\hookrightarrow Y\) is a cofibration: the finite
toric orbit-stratum decomposition can be triangulated with its
boundary spheres and edge arcs as a subcomplex. Thus this pushout
is a homotopy pushout and its Mayer--Vietoris and van Kampen
sequences apply. Van Kampen gives
\[
 \pi_1(H)=\langle a,b\mid ab^{-1}a^{-1}b=1\rangle
            =\mathbb Z^2.
\]
The preceding homology maps give
\[
 H_3(H)=\mathbb Z,\qquad
 0\longrightarrow\mathbb Z\longrightarrow H_2(H)
      \longrightarrow\mathbb Z\longrightarrow0,
 \qquad H_1(H)=\mathbb Z^2.
\]
The middle sequence splits as a sequence of abelian groups. Its
left generator is the image of \((0,[F])\) in the displayed
second-homology cokernel. It is therefore a primitive nonzero
class in \(H_2(H)\). There are no cells above dimension three.
This proves all assertions, including the nonbounding claim.
The endpoint equivalence classes, incidence ranks, commutator
abelianization, and primitive matrix are checked in
\path{code/d_push_cusp_orbit_cells.py}; the orientation argument
and comparison with the geometric pushout are supplied above.
\end{proof}
\begin{lemma}[C.3.8 (Proved): equivariant retraction of the cusp body]
For a sufficiently small closed cusp disc \(\overline\Delta\),
the compact manifold with boundary
\(V=f^{-1}(\overline\Delta)\) has a \(\rho\)-equivariant strong
deformation retraction onto \(W=f^{-1}(0)\). Consequently
\(V/\rho\) strongly deformation retracts onto \(H=W/\rho\),
and the fixed sphere \(F\) is primitive and nonzero in
\(H_2(V/\rho;\mathbb Z)\). It cannot bound a three-chain,
and hence cannot bound a three-ball, inside this cusp orbit body.
\end{lemma}
\begin{proof}
Write \(q\) for the cusp coordinate and set \(g=|q\circ f|^2\).
Choose a smooth invariant Hermitian metric by averaging any smooth
Hermitian metric over the compact circle. The negative gradient
flow of \(g\) is equivariant, fixes \(W\) pointwise, and points
strictly inward on \(\partial V\). Its flow \(\varphi_t\) exists
for every \(t\geq0\), since \(V\) is compact.

We prove the uniform convergence to \(W\) that is needed for a
strong deformation retraction. At every point of \(W\), the
actual semistable cusp charts give holomorphic coordinates with
\(q\circ f=z_1\cdots z_k\), where \(1\leq k\leq3\).
An initial holomorphic unit can be absorbed in one coordinate.
For a point off \(W\), some \(i\) satisfies
\( |z_i|\leq|q\circ f|^{1/k}\), and therefore
\[
 \left|\frac{\partial(q\circ f)}{\partial z_i}\right|
       \geq |q\circ f|^{1-1/k}.
\]
Uniform comparison with the chosen metric on a finite cover of
the compact fibre, and \(g\leq1\) after shrinking, give a
constant \(c>0\) such that
\[
 |\nabla g|\geq c g^{5/6}
\]
near \(W\). On the compact complement of a smaller such
neighbourhood, \(g>0\) and \(\nabla g\ne0\), because
\(f\) is a submersion off its central fibre. Decreasing \(c\)
extends this inequality to all of \(V\setminus W\).

Along a trajectory put \(h(t)=g(\varphi_t(x))\). Then
\(h'=-|\nabla g|^2\leq-c^2h^{5/3}\). If
\(G=\max_V g>0\), integration gives the uniform bound
\[
 h(t)\leq\bigl(G^{-2/3}+(2c^2/3)t\bigr)^{-3/2}.
\]
The remaining length of the trajectory starting at time \(t\)
is bounded by
\[
 \int_t^\infty|\nabla g(\varphi_s(x))|\,ds
 =\int_0^{h(t)}\frac{dv}{|\nabla g|}
 \leq\frac6c h(t)^{1/6}.
\]
One may first integrate to a finite time; the displayed uniform
bound and monotone convergence justify the endpoint expression.
Thus every trajectory has a limit in \(W\), and the tail-length
bound tends uniformly to zero. The maps \(\varphi_t\) converge
uniformly to a continuous equivariant retraction \(r:V\to W\).
The homotopy given by \(\varphi_{s/(1-s)}\) for \(s<1\)
and by \(r\) at \(s=1\) is continuous, equivariant, and fixes
\(W\) throughout. This is the required strong deformation
retraction. It descends to the orbit quotient, so the homology
and primitivity statement transfers from C.3.7. The same argument
does not supply a global Seifert surface in the full orbit sphere.
\end{proof}
\begin{lemma}[C.3.9 (Proved): the primitive circle marking]
Use the conductor generators \(a,b\) of C.3.7, with
\(E_1\) represented by the edge of direction \(e_1-e_2\),
\(E_2\) by the edge of direction \(e_1\), and the closing path
\(c\) on the fixed edge of direction \(e_2\). Under the
natural period marking of the cusp orbit fibre,
\[
 a\longmapsto\widehat u,\qquad
 b\longmapsto\widehat\gamma+\widehat u.
\]
In particular the circle obtained by gluing the two fixed boundary
spheres of the normalization has oriented generator
\(ab^{-1}=-\widehat\gamma\). It is primitive and has no
additional \(\widehat u\) component.
\end{lemma}
\begin{proof}
The cusp period action has fan translation
\(\lambda\mapsto B_0\lambda\), with
\[
 B_0=\begin{pmatrix}0&1\\-1&0\end{pmatrix},
 \quad \lambda=(\gamma,u)^t.
\]
This is the actual period matrix
\(Z=sB_0+C(q)\) of the cusp construction, not an unmarked
torus identification. The unit factors \(\exp(2\pi i C(q)\lambda)\)
fix the toric strata and hence do not change their triangle labels.
Take the cyclic fan directions
\[
 (D_1,\ldots,D_6)=
 (e_1-e_2,e_1,e_2,e_2-e_1,-e_1,-e_2).
\]
For \(L_v=(v,v+e_1,v+e_2)\) and
\(U_v=(v+e_1,v+e_2,v+e_1+e_2)\), the vertices of the
normalizing component have the following actual lifts:
\[
 \begin{array}{c|cccccc}
 &p_1&p_2&p_3&p_4&p_5&p_6\\\hline
 \text{triangle}&U_{-e_2}&L_0&U_{-e_1}&L_{-e_1}
                 &U_{-e_1-e_2}&L_{-e_2}.
 \end{array}
\]
For example \(p_1\) is the triangle
\(\{0,e_1-e_2,e_1\}=U_{-e_2}\); the other entries follow
by the same explicit substitution.

Choose \(A=L_0\) and the lift of \(c\) from
\(p_2=L_0\) to \(p_3=U_{-e_1}\).
The edge \(D_1\), oriented from \(A\) to \(B\), goes
from \(L_{-e_2}\) to \(U_{-e_2}\). Translate it by
\(+e_2\) to start at the selected lift \(L_0\).
It ends at \(U_0\), which is the chosen \(B\)-lift
translated by \(+e_1\). Thus \(a=E_1c^{-1}\) has
fan displacement \(+e_1=B_0\widehat u\).
The edge \(D_2\) oriented from \(A\) to \(B\)
starts already at \(L_0\) and ends at \(U_{-e_2}\).
Its displacement relative to the chosen \(B\)-lift is
\(e_1-e_2=B_0(\widehat\gamma+\widehat u)\).
This proves the two labels.

The real normalization circle is the path along \(D_1\)
and back along \(D_2\), closed by the gluing of the fixed
boundary spheres. Its image is \(E_1E_2^{-1}=ab^{-1}\),
so its label is \(-\widehat\gamma\). Equivalently, the
fixed-edge gluing has fan translation
\(B_0(1,0)^t=-e_2\), corresponding to the period
\(\widehat\gamma\); traversing its fundamental
domain in the stated direction has the opposite sign.
The two computed labels have determinant \(-1\), and hence
give the entire period lattice. A loop lift is legitimate even
at a vertex: a nonzero fan translation cannot preserve the finite
set of component labels incident to that vertex, so the period
action on the toric cover, also after the circle quotient, is free.
The finite triangle and displacement calculation is checked in
\path{code/d_push_cusp_marking.py}.
\end{proof}

\begin{proposition}[C.3.10 (Proved): reduction to a specified annular orbit map]
Let $\mathcal Q=X_{c_0}/\rho$ be the orbit space for the circle of
C.3.1. Choose a small cusp disc with coordinate $q$, and a smooth
radial function $\chi$ supported in that disc, equal to one on a
smaller disc. On its punctured annulus, let
$v_\gamma=\operatorname{pr}_{\delta}\Pi_{c_0}\widehat\gamma$
be the projected real period. Define an orbit map $\psi$ by
\[
 \psi([z])=[z+\chi(q)v_\gamma(q)]
\]
on the cusp annulus and by the identity elsewhere. This is a
well-defined diffeomorphism, equal to the identity near the fixed
set. It has a smooth equivariant lift to $X_{c_0}$. Every such lift
has mapping class $6[\Phi]$, for the positive $w$-loop convention
of C.1.1 and the closing map of C.3.5.

The assertion concerns the support of the \emph{orbit map}.
It does not assert an identity-ended lift supported in the annulus.
In particular it does not evaluate $6[\Phi]$.
\end{proposition}
\begin{proof}
We first specify the quotient transport used in the comparison.
Take $c(r)=c_0-12r$, $0\leq r\leq1$. On an ordinary fibre,
the quotient by $\rho$ is obtained by projecting the additive
period coordinates to $\mathbb C\times\mathbb R$, discarding
$\operatorname{Re}\zeta_2$. The period change
$(c(r)-c_0)e\otimes\gamma$ has real coefficient, so its
projected lattice is unchanged. This gives constant quotient
charts over the ordinary base.

The finite root quotient charts are also constant: the change of
their period columns and affine deck translations is in the real
$e$-direction and disappears under this projection. Their overlap
with the ordinary chart changes, however. With the notation and
signs of C.3.5,
\[
 \sigma_{j,c(r)}-\sigma_{j,c_0}
       =-12r\,\epsilon_j w_j e,
 \qquad w_j=(2\pi i)^{-1}\log s_j,
 \quad (\epsilon_1,\epsilon_2)=(1,-1).
\]
Let $R_a$ denote the positive real subgroup of the global
$\mathbb C^*$-action, with multiplier $e^a$. The transport in
the root chart which agrees with the ordinary identity is thus
$R_{12r\epsilon_j\log|s_j|}$ on the overlap. Choose a smooth
rotation-invariant real function $f_j(s_j)$ which vanishes near
zero and equals $\log|s_j|$ on that overlap. Replacing this by
$R_{12r\epsilon_j f_j(s_j)}$ extends the transport over the
root centre and commutes with its finite deck group. These
functions are smooth on the free root-cover quotient; no assertion
that $|t|^{1/m_j}$ is a smooth coordinate on the coarse base is
used. At the cusp, the varying multiplicative period unit is a
$\delta$-phase, so the toric quotient charts and their overlaps
are constant. We take the identity transport there.

This quotient transport is induced by a smooth equivariant
transport on the actual family. Here is the lifting justification.
On an ordinary or finite root chart the affine family and its
smooth lattice coordinates give local smooth horizontal lifts of
the prescribed quotient vector field. At the cusp use its smooth
toric charts, with the finite base parameter interval retained;
the dependence just removed is a circle phase, so these charts
give local equivariant lifts too, including at fixed points.
Two such lifts have the same quotient field and differ by a
vertical circle field on the free stratum. The difference vanishes
at a fixed point. An invariant smooth partition of unity on the
actual total space patches these local fields. Their parameter
component remains one and their quotient is the specified field.
Properness gives a flow for the whole compact interval. Denote
its endpoint by $H$. This argument uses smooth local lifts in the
given slices, rather than assuming an arbitrary orbit-space
isotopy lifts across the fixed set.

Put $G=\widetilde F_{12,c_0-12}\circ H$. By C.1.1 this smooth
equivariant diffeomorphism represents $6[\Phi]$. On ordinary
quotient charts, C.3.5 gives the positive real flow
$R_{a(t)}$, where
\[
 a(t)=\log|K(t)|,\qquad K(t)=t^4/(t-1)^3.
\]
On the finite root chart the same map is the flow with time
\[
 a_j(s_j)=\log|K(t)s_j^{-12\epsilon_j}|
                       +12\epsilon_j f_j(s_j).
\]
The first term is the logarithm of a nonvanishing unit and the
second is smooth; hence this time is smooth through the root
centre. On its outer overlap it equals $a(t)$. At the cusp the
quotient map is the bounded toric unit $U=U_{12,c_0-12}$ of
C.3.5. Write $u(q)$ for its additive projected logarithm.
That logarithm is single-valued on the full disc. The exact
period identity of C.3.5 is
\[
 u(q)+v_\gamma(q)=R_{a(q)}
\]
in the additive universal covering vector space; on the right
we mean the additive displacement of that real flow.

Multiplication by $U^{-s\chi(q)}$, $0\leq s\leq1$, defines a
smooth equivariant isotopy supported in the cusp. To see this
through every toric stratum, choose the explicit holomorphic
additive logarithm of the two nonzero unit coordinates of $U$.
Exponentiating $-s\chi$ times that logarithm gives a torus
element depending smoothly on $q$; multiplication acts smoothly
in every toric chart and commutes with the period identifications.
It is the identity near the outer edge of the disc.

There is a second global equivariant isotopy, the real
$\mathbb C^*$-flow with time $-s(1-\chi)a$, using $a_j$ in
the finite charts. The times agree on overlaps, are smooth on
the actual root quotients, and vanish on an inner cusp disc.
They are constant along the $\mathbb C^*$-orbits, so the inverse
flow is obtained by negating the time. Consequently this is an
isotopy of diffeomorphisms, not just a pointwise multiplication.

Compose $G$ with the endpoints of these two isotopies. Away
from the cusp its orbit map becomes the identity. In the cusp
annulus its projected displacement is exactly
\[
 a-\chi u-(1-\chi)a=\chi(a-u)=\chi v_\gamma.
\]
Where $\chi=1$, this is a whole projected lattice period and
the orbit map is again the identity; on the central fibre the
cusp-unit description gives the same identity directly. Thus the
result is an actual smooth equivariant lift of the stated $\psi$,
and has the same mapping class as $G$.

Finally $A_0\widehat\gamma=\widehat\gamma-\widehat\delta$
makes $v_\gamma$ single-valued on the punctured cusp after the
$\delta$-quotient. Translation by $\chi v_\gamma$ and its
negative are smooth mutually inverse maps on that free orbit
annulus, and they are the identity on collars of its ends.
This proves the direct well-definedness and smoothness assertion.
C.3.6 shows that every other existing smooth equivariant lift has
the same mapping class. This construction supplies no assertion
that the full lift is supported in the annulus: its residual
circle gauge outside that annulus has not been discarded.
\end{proof}
\begin{proposition}[C.3.11 (Proved): a nonzero normal Euler number]
Glue the two fixed boundary spheres of \(Y=\mathrm{dP}_6/\rho\)
using their actual conductor identification, and call the resulting
closed oriented three-manifold \(N\). Its natural map to the
smooth cusp orbit five-manifold is an immersion. Its sole
self-intersection circle has two disjoint preimage circles, the
two pairs of nonfixed edge intervals of C.3.7.
Along the fixed sphere \(F\subset N\), the oriented normal
two-plane bundle of this immersion has Euler number of absolute
value two. If its complex orientation is that of the Hopf invariant
\(a=xz\), where \(x,z\) have circle weights \(+1,-1\), the
Euler number is \(-2\).
\end{proposition}
\begin{proof}
Away from the conductor the normalization quotient is embedded.
At an ordinary nonfixed double point the circle acts freely along
the double curve, and the two smooth three-dimensional branches
intersect transversely along its one-dimensional orbit quotient.
At a fixed double point the two normal coordinates have weights
\(+1,-1\). The Hopf quotient coordinates are
\[
 a=xz\in\mathbb C,\qquad
 b=(|x|^2-|z|^2)/2\in\mathbb R,
\]
together with the coordinate along the fixed curve. The two
halves of \(Y\) glue as the two signs of \(b\), so their
union is smooth. At a triple point the full base equation is
\(q=xyz\), with \(y\) tangent to \(F\) and invariant.
In the quotient it becomes \(q=ay\). Hence the central orbit
fibre is the union of the two smooth three-planes
\(\{a=0\}\) and \(\{y=0\}\), intersecting transversely
along the \(b\)-axis. This verifies immersion and transversality
at both previously distinguished vertices. The endpoint
identifications of C.3.7 join the two nonfixed double arcs into
one circle; its preimages in \(N\) are disjoint circles.

The two complex weight lines in \(N_{F/X}\) are each
\(\mathcal O_{\mathbb P^1}(-1)\). To see this with the actual
normalization, each is the normal line of one fixed edge in its
\(\mathrm{dP}_6\) branch, and each edge has self-intersection
\(-1\). This remains true at the two vertices because those
branches and their edge normals are smooth there. The circle
characters distinguish the lines globally, so no choice of a
splitting of an extension is required.

The Hopf invariant \(a=xz\) is valued in the tensor product
of these two weight lines. It therefore gives the complex normal
line of the branch \(\{a=0\}\) along \(F\) as
\[
 \mathcal O(-1)\otimes\mathcal O(-1)=\mathcal O(-2).
\]
At a triple point \(q=ay\), the coordinate \(a\) still is
a normal coordinate of that branch. The zero of \(y\) does
not alter its normal bundle; using \(q\) instead would lose
precisely this distinction. The stated Euler number is consequently
\(-2\) for the \(a\)-complex orientation and \(+2\) if that
normal orientation is reversed.

In particular this immersion cannot simply be compressed along
\(F\) to an immersed three-sphere while retaining its normal
two-plane bundle: extension over a bounding three-disc would force
its Euler pairing on the boundary sphere to vanish. This identifies
a required framing correction; it does not compute that correction
or a Haefliger invariant for the annular twist.
\end{proof}
\begin{proposition}[C.3.12 (Proved): obstruction to the literal translational Seifert quotient]
Fix an admissible parameter and a small closed cusp body $V$, with
central fibre $W$. On every regular fibre let $K_q$ be the real
three-dimensional subtorus marked by $\ker\gamma$, namely by
$\widehat u,\widehat w,\widehat\delta$.
If $Z$ is Hausdorff and a continuous map $h:V\to Z$ is constant
on every translate of $K_q$ for $q\ne0$, then $h|_W$ is constant.

Consequently there is no continuous extension of this literal
regular-fibre translational quotient whose central fibre is a
circle. In particular, the translational quotient in the first
paragraph of \citet[Remark 8.3]{E26} cannot be identified with the
circle-valued cusp moment quotient described there while retaining
those equivalence classes on every regular fibre. A construction
which changes the regular-fibre equivalence relation is not excluded.
This correction concerns the rank-three quotient; it does not affect
the smooth $\widehat\delta$-circle quotient in C.3.10.
\end{proposition}
The nonextension of the raw invariant-circle coordinate was already
recorded in the supplied \path{background/notes/A1_notes.md},
Section~2, item A1--1, and
\path{background/S6_SURVEY_20261112.md}, item F9.
The assertion proved here has the stronger quantified scope of every
Hausdorff-valued continuous map with the stated regular translational
leaves. No claim of priority is made.
\begin{proof}
We use the actual cusp toric coordinates and periods. Write
\[
 q=e^{2\pi i s},\qquad
 Z=sB_0+C(q),\qquad
 B_0=\begin{pmatrix}0&1\\-1&0\end{pmatrix},
 \qquad C_{22}(q)=\mu(q).
\]
The columns correspond to $(\widehat\gamma,\widehat u)$;
the unitary torus directions are $(\widehat w,\widehat\delta)$.
On the dense torus, with multiplicative coordinates $(z_1,z_2)$,
translation by the real multiple $r\widehat u$ is exactly
\begin{equation}\label{eq:C12-u-flow}
 (z_1,z_2)\longmapsto
 \bigl(q^r e^{2\pi i r C_{12}(q)}z_1,
             e^{2\pi i r\mu(q)}z_2\bigr).
\end{equation}
Here a branch is chosen on the positive real $q$-ray. This is a
translation in $K_q$ for every real $r$, not merely for integral $r$.
The two remaining circle directions rotate $z_1$ and $z_2$
independently.

The limiting constant $\mu_0=\mu(0)$ has nonzero imaginary part.
For the branch used in D.4.7 it is
\[
 \mu_0=\frac12-\frac{i}{2\pi}\log(5+2\sqrt6),\qquad
 |e^{2\pi i\mu_0}|=5+2\sqrt6>1.
\]
Only the nonvanishing of its imaginary part is needed. It can
also be checked directly without that branch choice: the cusp
limit of $\wp(\mu;\tau)=0$ is
$\pi^2\csc^2(\pi\mu_0)-\pi^2/3=0$, so
$\sin^2(\pi\mu_0)=3$, which is impossible for real $\mu_0$.
The Fourier limit and holomorphic finite cusp branch are the
ones checked in D.4.7.

Consider the edge between the lattice vertices $0$ and $e_1$
of the Mumford fan. On the open part of its two adjacent
components the toric chart is
\[
 x=z_1,\qquad y=q/z_1,\qquad z=z_2\in\C^*,
 \qquad xy=q.
\]
This chart follows directly from the two primitive height-one
rays $(0,0,1)$ and $(1,0,1)$; its additional invertible character
is $z_2$. It is also the ordinary double-curve chart of the
fixed bottom triangulation. Let $P(x,z)$, for $x,z\in\C^*$,
denote its central point $(x,0,z)$, and let $Q(z)$ denote
$(0,0,z)$ on the double curve. We apply these notations also
to their images in the period quotient. The set of all $P(x,z)$
is the dense torus of the normalizing $\mathrm{dP}_6$ component,
whose image is dense in the irreducible central fibre $W$.

First $h(P(x,z))$ is independent of $x$. To check positive
radial changes, fix $\lambda\in\R$, put $q_n=e^{-n}$, and
translate $(x,q_n/x,z)$ by
$r_n\widehat u$, where $r_n=-\lambda/n$. Since $C_{12}$ and
$\mu$ are bounded, the two multiplicative factors in
\eqref{eq:C12-u-flow} tend respectively to $e^\lambda$ and
$1$. The translated sequence therefore tends to $P(e^\lambda x,z)$,
whereas the original sequence tends to $P(x,z)$. Their $h$-images
agree term by term. Continuity and the Hausdorff property give
equality of the two limiting images. Rotation by
$\widehat w$ treats the argument of $x$ in the same way.
Thus we may write $H(z)=h(P(1,z))$.

Next fix $0<r<1$ in \eqref{eq:C12-u-flow} and let $q\downarrow0$
on the positive real ray. Starting at $(x,q/x,z)$, the translated
coordinates satisfy
\[
 x_r=q^r e^{2\pi i r C_{12}(q)}x\longrightarrow0,\qquad
 y_r=q^{1-r} e^{-2\pi i r C_{12}(q)}/x\longrightarrow0,
\]
and $z_r\longrightarrow e^{2\pi i r\mu_0}z$. Hence the same
continuity argument gives the exact equality
\begin{equation}\label{eq:C12-limiting-identification}
 H(z)=h\bigl(Q(e^{2\pi i r\mu_0}z)\bigr)
 \qquad(0<r<1).
\end{equation}
For $r_1,r_2\in(0,1)$ apply this equality first to $z$ and
then to $z'=e^{2\pi i(r_1-r_2)\mu_0}z$. Their double-curve
limits are the same point, so
\[
 H(z)=H(e^{2\pi i a\mu_0}z)\qquad(-1<a<1).
\]
Iteration gives the equality for every real $a$. Rotation by
$\widehat\delta$ also gives $H(z)=H(e^{i\vartheta}z)$ for every
$\vartheta$. Since $\operatorname{Im}\mu_0\ne0$, the moduli
$|e^{2\pi i a\mu_0}|$ range over all positive real numbers.
Combining these two invariances proves that $H$ is constant on
$\C^*$. Therefore $h$ is constant on the dense torus in $W$.
The inverse image of a point of the Hausdorff space $Z$ is
closed; density proves constancy on all of $W$.

For comparison, the raw real period coordinate already fails
continuity in an elementary fixed-edge chart. Take
$q=e^{-2\pi T}$, $z_1=1$, $z_2=q^r$, with $0<r<1$ and
$T\to+\infty$. In the chart
$(z_2,q/z_2,z_1)$ all these paths tend to the same point
$(0,0,1)$. Their logarithmic modulus vectors are $(0,rT)$.
Since $\operatorname{Im}Z=TB_0+O(1)$, their real period
vectors are
\[
 (\operatorname{Im}Z)^{-1}(0,rT)^t
                  =(-r,0)^t+O(T^{-1}).
\]
Thus their raw $\gamma$ circle coordinates tend to $-r$ modulo
$\Z$, with different limits at that one central point. The
stronger argument \eqref{eq:C12-limiting-identification} shows
why an arbitrary base-dependent reparametrization of those
circle coordinates cannot repair the asserted central circle
while keeping the same regular-fibre translational leaves.
The exact limiting multiplier and the inverse-matrix calculation
are independently checked in
\path{code/d_push_cusp_quotient_obstruction.py}; the continuity
and closure argument above, rather than that finite check,
proves the universal assertion.
\end{proof}
\begin{proposition}[C.3.13 (Proved): a corrected marked Seifert comparison]
There is a smooth map $\mathcal G:X_c\to S^3$, invariant under
the $\widehat\delta$ circle, with the following properties.
It covers the given base coordinate through the explicit weighted
Seifert map
\[
 p:S^3\longrightarrow\PP^1,\qquad
 p(z_1,z_2)=\frac{z_2^3}{z_2^3-z_1^4},\qquad
 \lambda\cdot(z_1,z_2)=(\lambda^3z_1,\lambda^4z_2).
\]
Outside a chosen small cusp body it is the actual quotient by
the regular $\ker\gamma$ subtori, followed by a specified
base-dependent circle gauge and the explicit identification below.
It includes the actual finite root quotients with rotation data
$+1/3$ and $-1/4$. Inside the cusp body the map is modified;
it does not preserve the original translational leaves.

On the outer cusp collar a positively marked $\widehat\gamma$
circle maps orientation-preservingly onto a regular weighted
Seifert fibre. The central fibre of the target is the positive
$(3,4)$ torus knot, and a nearby regular fibre has linking number
$12$ with it, in the standard orientation of $S^3\subset\C^2$.
This is framing data in the constructed three-dimensional target.
No identification of its normal framing with the six-dimensional
orbit-knot framing, and no value of $6[\Phi]$, follows from this
statement.
\end{proposition}
\begin{proof}
We first fix the real character and all of its overlap corrections.
On an ordinary fibre, let
$\gamma_z:\C^2\to\R$ be the real-linear functional determined by
\[
 \gamma_z(\Pi_c(z)v)=\gamma(v)\qquad(v\in\R^4).
\]
The real period matrix is invertible in the admissible range.
The resulting circle coordinate
$\theta=\gamma_z(\zeta)\pmod\Z$ is invariant under ordinary
monodromy because $\gamma A_j=\gamma$. It is the literal
regular-fibre quotient character, not a new choice of marking.

On the $j$th finite root chart write $\theta_j$ for this
character before the affine finite quotient. The exact map from
that chart to the ordinary chart is translation by
$\sigma_{j,c}=w_j\Pi_cv_j$, where
$w_j=(2\pi i)^{-1}\log s_j$ and
$\epsilon_j=\gamma(v_j)$ is $1,-1$, respectively.
Consequently its real circle overlap is exactly
\begin{equation}\label{eq:C13-real-overlap}
 \theta=\theta_j+\epsilon_j\operatorname{Re}w_j+a_j^{\mathrm{rad}}(t),
 \qquad
 a_j^{\mathrm{rad}}(t)=\operatorname{Im}w_j\,
                  \gamma_z(i\Pi_cv_j).
\end{equation}
The last expression is real and single-valued on the punctured
base disc. Indeed $\operatorname{Im}w_j$ is single-valued,
$v_j$ is fixed by the local integral monodromy, and its
complex-linear covariance preserves $\gamma_z$.
No logarithmic radial term has been dropped.

Choose disjoint finite discs and smooth radial cutoffs equal to
one on smaller discs and zero outside the larger ones. They define
a real smooth function $d(t)$ on $B^\circ$ which is $-a_j^{\mathrm{rad}}(t)$
on the smaller punctured discs and is zero on a cusp collar.
Choose the roots $F^3=t$, $G^4=t-1$ on the uniformizing cover and
the logarithm $\ell=(2\pi i)^{-1}\log(F/G)$ used in C.3.5.
Put
\[
 \Theta=\theta-\operatorname{Re}\ell+d(t).
\]
Let $R_t>0$ be the unique solution of
\[
 R_t^6|G|^2+R_t^8|F|^2=1.
\]
Existence and uniqueness follow from strict increase from zero
to infinity; the derivative in $R_t$ is positive, so the
implicit-function theorem gives the required smooth dependence.
The proposed map on the ordinary region is
\begin{equation}\label{eq:C13-weighted-map}
 (z_1,z_2)=
 \bigl(R_t^3G e^{6\pi i\Theta},\,
                R_t^4F e^{8\pi i\Theta}\bigr).
\end{equation}
Its norm is one and direct substitution gives $p(z_1,z_2)=t$.
Increasing $\theta$ by $a$ is exactly the weighted action of
$e^{2\pi ia}$, so the positive $\gamma$ orientation is retained.

We check descent and the finite centres explicitly. For the
clockwise first root generator $F$ changes by $e^{-2\pi i/3}$,
$G$ is unchanged, and $\ell$ changes by $-1/3$.
Thus $\Theta$ changes by $+1/3$, and both coordinates in
\eqref{eq:C13-weighted-map} remain unchanged. For the second
generator $G$ changes by $e^{-2\pi i/4}$, $F$ is unchanged,
and $\Theta$ changes by $-1/4$. Again both coordinates are
unchanged: their total exponents are respectively $-1$ and
$-1$. These generators and changes of logarithm prove descent
on the entire ordinary base.

On a full finite root disc the exact identity from C.3.5 is
$\ell=\epsilon_jw_j+\kappa_j$, where $\kappa_j$ is a
holomorphic logarithm of an invariant nonvanishing unit.
Equations \eqref{eq:C13-real-overlap} and $d=-a_j^{\mathrm{rad}}$ therefore give
\[
 \Theta=\theta_j-\operatorname{Re}\kappa_j.
\]
This expression is smooth through the root centre. The vanishing
factor $F$ at the first centre, or $G$ at the second, is the
root coordinate times a holomorphic unit. The other factor is
nonzero. Thus \eqref{eq:C13-weighted-map} extends smoothly to
the full root chart. Its transverse derivative is nonzero and
the circle factor has degree one; it identifies the actual
finite quotient solid torus with the corresponding weighted
solid torus. The deck actions on $(s_j,\theta_j)$ are precisely
\[
 (s_1,\theta_1)\mapsto(e^{-2\pi i/3}s_1,\theta_1+1/3),
 \qquad
 (s_2,\theta_2)\mapsto(e^{-2\pi i/4}s_2,\theta_2-1/4).
\]
This verifies the overlap maps, not merely the exceptional
multiplicities of an abstract Seifert space.

We next verify the cusp meridian and construct the replacement
map upstairs, without assuming a smooth quotient structure at
the fixed set. Let $V$ be a small smooth compact cusp body and
$\overline V=V/\rho$. C.3.7--C.3.9 identify
$H_1(\overline V;\Z)=\Z\widehat\gamma\oplus\Z\widehat u$.
The outer orbit link is the mapping torus of the action
$u\mapsto u+w$ on the $(u,w)$ torus, times the $\gamma$ circle.
Its first homology is therefore freely generated by
$\widehat\gamma,\widehat u$ and a transverse cusp meridian;
$\widehat w$ is a commutator. The two first generators have
the stated images in $H_1(\overline V)$ by the actual fan marking
in C.3.9.

The transverse meridian has zero $\gamma$ value and bounds
an actual disc in $V$. To see this without choosing a meridian
up to an unknown fibre twist, take the toric section
$z_1=z_2=1$. In the chart of C.3.12 it is
$(x,y,z)=(1,q,1)$, so it extends across $q=0$.
Its additive logarithms are zero and hence $\theta=0$
identically on its boundary. Thus the integral class on
$\overline V$ with values $1,0$ on
$(\widehat\gamma,\widehat u)$ restricts to exactly the raw
character $[d\theta]$ on the orbit link, including the transverse
meridian value. Pull it back to $V$.

Represent the pulled-back integral class by a smooth closed real
one-form $\alpha$ with integral periods. Average it over the
$\delta$ circle. Each group element is isotopic to the identity,
so the cohomology class and its integral periods are unchanged.
The averaged form is invariant. If $K$ denotes the circle
generator, Cartan's formula gives $d(\alpha(K))=0$.
The body is connected and has fixed points, where $K=0$;
therefore $\alpha(K)=0$ everywhere.
On an outer collar $\alpha-d\theta$ is exact, since the two
forms have the same periods. Choose its real primitive,
average the primitive over the circle, and extend it to $V$
using an invariant collar cutoff. Subtracting its differential
from $\alpha$ makes the resulting invariant integral closed
form equal to $d\theta$ on a smaller outer collar. Integrating
it modulo $\Z$, with the constant chosen at one collar point,
gives a smooth map
\[
 g:V\longrightarrow\R/\Z
\]
which equals $\theta$ on that collar. Its contraction with $K$
is zero, so $g$ is invariant. This construction is entirely on
the smooth manifold $V$ and yields a continuous factor on
$\overline V$.

For completeness, the prescribed orbit boundary map has a
unique relative homotopy class of extensions to $\overline V$.
The restriction of $H^1(\overline V;\Z)$ to the link is
injective by the two period generators just checked, and both
spaces are connected. The relative cohomology exact sequence
therefore gives $H^1(\overline V,\partial\overline V;\Z)=0$.
Maps to $\R/\Z$ with fixed boundary values are classified by
this relative group: a map represents its integral period class,
and a map with zero class has a real lift which supplies the
relative homotopy. This assertion does not classify normal
framings of any embedded trace.

Use now the holomorphic cusp coordinate $q=1/t$; changing from
the original Mumford coordinate multiplies it by a unit and
does not change the meridian just verified. Write
$q=e^{2\pi i s}$. Compatible roots satisfy
\[
 F=q^{-1/3},\qquad G=q^{-1/4}(1-q)^{1/4},\qquad
 \ell=-s/12+\kappa_0(q),\qquad
 \kappa_0(q)=-\frac{\log(1-q)}{8\pi i}.
\]
A different compatible root choice only changes the displayed
unit gauge. The smooth fibre coordinate in the weighted model
near its regular cusp knot is
\begin{equation}\label{eq:C13-cusp-gauge}
 \vartheta=\Theta-\operatorname{Re}s/12
           =\theta-\operatorname{Re}\kappa_0(q),
\end{equation}
because $d=0$ there. In particular no additional integer
meridian winding is present. This is the required boundary
framing comparison; the finite labels alone would not give it.

Let $S(q)>0$ solve
$S(q)^6|1-q|^{1/2}+S(q)^8=1$. The explicit smooth solid-torus
chart of the target is
\[
 \Phi_0(q,\vartheta)=
 \bigl(S(q)^3(1-q)^{1/4}e^{6\pi i\vartheta},
                         S(q)^4e^{8\pi i\vartheta}\bigr).
\]
It is a diffeomorphism onto a neighbourhood of the regular
cusp knot. Indeed its inverse has
$q=1-z_1^4/z_2^3$ and
$e^{2\pi i\vartheta}=z_2(1-q)^{1/4}/(S(q)z_1)$;
all factors are nonzero there. Define the replacement on $V$
by
\[
 \mathcal G(x)=
 \Phi_0\bigl(q(x),g(x)-\operatorname{Re}\kappa_0(q(x))\bigr).
\]
It agrees exactly with \eqref{eq:C13-weighted-map} on the outer
collar by \eqref{eq:C13-cusp-gauge}, and hence glues with the
ordinary and finite maps already constructed. Smoothness and
$\delta$ invariance hold upstairs on every piece. The first
coordinate of the solid-torus chart proves that it still covers
$f$. By C.3.12 its new cusp leaves cannot be the original
translational leaves on every punctured fibre; that change is
an explicit part of the construction.

Finally the knot $q=0$ in the target is
$z_2^3=z_1^4$, parametrized with positive orientation by the
weighted circle $(e^{3i\phi},e^{4i\phi})$ with the fixed positive
radii prescribed by $S(0)$. It is the positive $(3,4)$ torus
knot. For a small fixed $q\ne0$, along the nearby fibre
\[
 (z_2^3-z_1^4)\circ\Phi_0(q,\vartheta)
                  =q S(q)^{12}e^{24\pi i\vartheta}.
\]
The argument therefore has degree $12$. The function
$z_2^3-z_1^4$ has a simple complex-normal zero along the
knot in $S^3$, so its argument represents the positive
meridian class in the knot complement. This degree is its
linking number with the knot. The orientation is the one
already matched to positive $\gamma$ in
\eqref{eq:C13-weighted-map}. The computation establishes the
three-dimensional push-off datum only. The actual normal-line
comparison and its contribution are established separately in
C.3.20--C.3.22; they are not inferred from this linking number.
The finite phases, radial cancellation, link abelianization and
polynomial identities are checked exactly in
\path{code/d_push_corrected_seifert.py}.
\end{proof}
\begin{proposition}[C.3.14 (Proved): an actual characteristic-number target]
The smooth orbit pair for the circle of C.3.1 is diffeomorphic to
the standard pair $(S^5,S^2)$. Its free complement, after deleting
a small fixed tubular neighbourhood, is therefore
$P\simeq S^2\times D^3$. The circle action on $X_c$ is smoothly
equivariantly equivalent, with its circle direction fixed, to the
linear Hopf action on
$\partial(D^4_{\mathbb H}\times D^3)$.

Choose such an identification and regard the actual annular map
$\psi$ of C.3.10 as a diffeomorphism of $P$ equal to the identity
near its boundary. Let $\widetilde\psi$ be its extension by the
identity to the double $DP\simeq S^2\times S^3$. Form
\[
 M_\psi=(S^2\times D^4)\cup_{\widetilde\psi}
                  -(S^2\times D^4).
\]
This is a closed simply connected spin six-manifold with
$H^2(M_\psi;\Z)=H^4(M_\psi;\Z)=\Z$ and $H^3(M_\psi;\Z)=0$.
Let $x$ be the generator restricting to the oriented meridian
generator on each $S^2\times D^4$, and put
\[
 k_\psi=\langle x^3,[M_\psi]\rangle,\qquad
 J_\psi=\frac{\langle p_1(M_\psi)x-4x^3,[M_\psi]\rangle}{24}.
\]
Both are integers. The unframed knot coordinate of the orbit
knot of the clutching sphere for $6\Phi$ is $-J_\psi$ in
Wall's surgery convention. Consequently, with one fixed sign
depending on comparison of that convention with the generator
of C.1.1,
\[
 \mu(\Sigma_{6\Phi})=\pm\frac{J_\psi}{28}\pmod\Z.
\]
The pure framing subgroup has zero image in $\Theta_7$.
No value of $J_\psi$ or $k_\psi$ is asserted here.
\end{proposition}
\begin{proof}
The dimension hypothesis in
\citet[Section 1, Propositions 1.3--1.4]{FintushelPao1983D}
is $n>5$. Applied to this semifree action on $S^6$, it gives a
smooth orbit homotopy five-sphere and a smooth embedded fixed
two-sphere. Since $\Theta_5=0$ \citep{KervaireMilnor1963}, the
orbit is the standard $S^5$. The unknotting range
$C_n^q=0$ for $n<2q-3$
\citep[Corollary 6.6]{Haefliger1966D} applies at $(n,q)=(2,3)$.
Thus the fixed embedding is standard and its exterior is
$S^2\times D^3$.

Here the orbit identification can be lifted to the claimed
equivariant identification. In the normal bundle of the fixed
curve, conjugating the negative weight line converts the circle
representation to two scalar weight-one complex lines. By
C.3.11 their Chern classes are $-1$ and $+1$. The resulting
rank-two complex bundle over $S^2$ has first Chern class zero,
and is trivial, since clutching rank-two complex bundles over
$S^2$ is classified by $\pi_1(U(2))$ and its determinant.
Thus its equivariant tubular neighbourhood is the linear model.
The oriented normal three-plane bundle of the orbit sphere is
also trivial. The ambient pair diffeomorphism can be adjusted
on a tubular neighbourhood to this prescribed normal model:
the relevant change of oriented framing is a map
$S^2\to SO(3)$, and $\pi_2(SO(3))=0$.

On the complement, the principal circle bundle has Euler class
equal to a primitive generator of $H^2(P;\Z)$, as is seen on
each normal Hopf sphere. Choose the orbit identification to
match this signed generator with the linear model. Principal
circle bundles are classified by that class, so their complements
are equivariantly isomorphic. On the common boundary the
difference from the prescribed tubular isomorphism is a gauge
$S^2\times S^2\to S^1$. It has a real logarithm because
$H^1(S^2\times S^2;\Z)=0$. Extending that logarithm over a
collar and using a cutoff corrects the bundle isomorphism there.
The two isomorphisms now glue smoothly. This proof uses the
linear tubular model at the fixed set; it does not assume that
an arbitrary quotient diffeomorphism lifts smoothly there.

The support of $\psi$ misses the fixed neighbourhood by C.3.10.
Any boundary-fixed diffeomorphism of $P$ preserves the meridian
Euler class, since restriction $H^2(P)\to H^2(\partial P)$ is
injective. It therefore has a circle-bundle lift, which can be
made the identity on a boundary collar by the same logarithmic
gauge correction. Extend that lift by the identity on the fixed
neighbourhood. C.3.6 compares it with the lift actually constructed
in C.3.10, so its mapping class is $6[\Phi]$.

Clutch two linear seven-discs along this lift. Their fixed discs
clutch by the identity to a fixed $S^3$. Their orbit six-discs
clutch along $\psi$ on $S^5$, giving an orbit homotopy six-sphere,
which is standard because $\Theta_6=0$. Denote the resulting
orbit knot by $K_\psi\subset S^6$. This is the orbit knot of
the actual clutching sphere for $6\Phi$.

We spell out its framed surgery, so that the cap is specified.
In each orbit six-disc, the complement of a standard fixed
three-disc neighbourhood is $N=S^2\times D^4$. Its boundary
is divided into an outer and an inner copy of $P=S^2\times D^3$.
The two outer copies glue by $\psi$ to form the knot exterior.
The inner copies give its boundary $S^2\times S^3$ with the
framing inherited from the fixed-disc normal products. Framed
surgery fills that boundary by a third copy of $N$.
Combine this third copy with one of the first two along its
inner $P$. Two four-discs glued along a boundary three-disc
form a four-disc. The combined piece is therefore again $N$,
and the remaining gluing map is $\psi$ on the outer half and
the identity on the inner half. This is exactly the displayed
$M_\psi$, with the boundary orientation convention used there.

Van Kampen gives $\pi_1(M_\psi)=0$. Mayer--Vietoris for its
two copies of $S^2\times D^4$, with intersection $S^2\times S^3$,
gives the asserted cohomology groups and no torsion. The
restriction map in degree two is injective, and $w_2$ restricts
to zero on both pieces; thus $w_2(M_\psi)=0$. The spin
structure is unique. If $L_x$ is the line bundle with Chern
class $x$, the twisted spin index formula used in C.2.1 gives
\[
 \operatorname{ind}(D\otimes L_x)
 =\left\langle\frac{x^3}{6}-\frac{p_1x}{24},[M_\psi]\right\rangle
 =-J_\psi\in\Z.
\]
This also verifies the required integrality independently of the
classification of six-manifolds. The identical characteristic
integer is recorded in \citet[Section 3, page 1911]{Olbermann2010D}.

For a single framed three-knot, Wall's direct surgery calculation
gives
\[
 k=6\varphi+\beta',\qquad \langle p_1x,[M]\rangle=4\beta',
\]
where $\varphi$ is its unframed knot integer and $\beta'$ its
framing integer
\citep[Theorem 4, pages 360--361]{Wall1966VE}.
Hence $J=-\varphi$. The original pages were inspected directly;
the published corrigendum changes Theorem 14, not this formula
\citep{Wall1967CorrE}. Finally the orbit-knot classification and
\citet[Example 7.14 and Proposition 7.15]{Davis1982D} send a
knot generator to the Milnor sphere $M^7_3=S_{2,-1}(\mathbb H)$.
The associated oriented rank-four bundle over $S^4$ has Euler
evaluation one and first Pontryagin evaluation six. Thus C.2.1
gives $(6^2-4)/896=1/28$ with the corresponding
orientation. Comparing the two choices of knot generator gives
the fixed sign in the statement.

The framing kernel has a useful direct verification. Represent
its generator on $P$ by
$(a,y)\mapsto(r(y)a,y)$, with
$r:(D^3,\partial D^3)\to(SO(3),1)$. Lift $r$ to
$R:D^3\to SU(2)$, equal to one near the boundary. In the
quaternionic linear model its equivariant lift is
$(v,y)\mapsto(R(y)v,y)$ on
$\partial(D^4_{\mathbb H}\times D^3)$; the circle acts by
right complex multiplication. This formula extends over the
seven-disc. Its clutching sphere is therefore standard.
The framing direction cannot change the $\Theta_7$ class,
in agreement with the formula $J=-\varphi$.
\end{proof}

\begin{corollary}[C.3.15 (Proved): the specified cap integer vanishes]
The characteristic integer of this actual annular cap is
$J_\psi=0$ in $\Z$.
\end{corollary}
\begin{proof}
The complete evaluation is given below in C.3.18--C.3.22.
The actual affine-frame defect is $p_{\mathrm{rel}}=-24$,
the actual delta-line relative square is $s_\delta=-2$, and
the common signed cap formula is
\[
 J_\psi=\epsilon_{\rm cap}
              (p_{\mathrm{rel}}-12s_\delta)/24=0.
\]
This uses the global frame and Picard-line comparisons, not
only the local unipotent derivative or a target linking number.
The fixed sign comparing Wall's nonzero generator with the
positive Eells--Kuiper generator is immaterial for this zero.
C.3.23--C.3.24 give the integer-level argument for the original
one-loop monodromy.
\end{proof}
\begin{proposition}[C.3.16 (Proved): the actual cusp-link inclusion]
Orient the central orbit cycle $H=W/S^1_\delta\subset Q=X_c/S^1_\delta$
by the complex $q$-normal orientation, where $q=1/t$ near the cusp.
Use the marked positive period $\gamma$ and the actual transverse
$q$-meridian of C.3.13. Then
\[
 H_1(Q\setminus H;\Z)=\Z,
 \qquad \operatorname{lk}(\gamma,H)=12.
\]
More precisely, write the cusp orbit link as
$L=S^1_\gamma\times\operatorname{Nil}^3(u,w,q)$, with the integral
basis $(\gamma,u,q)$ of $H_1(L;\Z)$ fixed in C.3.13. The inclusion
homomorphism is
\[
 H_1(L;\Z)\longrightarrow H_1(Q\setminus H;\Z),
 \qquad (a,b,c)\longmapsto 12a+c,
\]
where the target generator is the positive meridian of $H$.
This is a statement about the actual cusp inclusion, not only the
weighted target three-sphere.
\end{proposition}
\begin{proof}
By C.3.14, $Q$ is the standard smooth five-sphere. The finite
normalization pushout in C.3.7 makes $H$ a compact, locally
contractible space, and gives $H^3(H;\Z)=\Z$. Its oriented
normalization has degree one at a generic smooth point; identification
along the conductor, of dimension at most two, does not change the
top-dimensional fundamental cycle. Thus $H$ has the multiplicity-one
oriented fundamental cycle specified in the statement. Integral
Alexander duality gives
\[
 \widetilde H_1(Q\setminus H;\Z)
 \simeq \widetilde H^3(H;\Z)=\Z.
\]
The local positive meridian at any generic smooth point pairs to one
with that fundamental cycle and therefore represents the indicated
primitive generator.

Let $\overline G:Q\to S^3$ be the continuous descent of the smooth
$S^1_\delta$-invariant map constructed in C.3.13, and write its two
complex coordinate functions as $G_1,G_2$. Define the globally
single-valued continuous complex function
\[
 h=G_2^3-G_1^4.
\]
Since the weighted target map is
$p(z_1,z_2)=z_2^3/(z_2^3-z_1^4)$ and
$p\circ\overline G=\overline f$, the zero set of $h$ is exactly $H$.
In the cusp solid-torus formula of C.3.13, it is explicitly
\[
 h=q\,S(q)^{12}\exp(24\pi i\vartheta),\qquad
 \vartheta=g-\operatorname{Re}\kappa_0(q),
\]
where $S(q)>0$, $g$ is circle-valued, and $\kappa_0$ is a
single-valued holomorphic function on the cusp disc. The factor
multiplying $q$ is a continuous nowhere-zero function. Hence on a
small positive normal circle at a generic smooth point of $H$, the
map $h/|h|$ has degree $+1$. It follows that
\[
 [h/|h|]\in H^1(Q\setminus H;\Z)
\]
is exactly the integral meridian, or linking, class. This argument
requires only continuity of the descended map at fixed points; it
does not assert that an arbitrary invariant smooth map descends
smoothly in every orbit-space coordinate system.

On the outer cusp collar the extension $g$ equals the raw marked
character $\theta$. Traversing the positive $\gamma$ period changes
$\theta$ by one, while $q$ and the other torus coordinates stay fixed.
The displayed formula for $h$ consequently has argument degree
$12$ on that loop. A traversal of $u$ changes neither $q$ nor
$\theta$, and has degree zero. Finally, the actual positive
transverse meridian is the section $z_1=z_2=1$ of C.3.13. On it
$\theta=0$ identically, so
\[
 \frac{h}{|h|}
 =\frac{q}{|q|}\exp\bigl(-24\pi i\operatorname{Re}\kappa_0(q)\bigr).
\]
The second factor extends over the whole disc and has degree zero.
The meridian therefore has degree one. These three evaluations give
the asserted inclusion homomorphism. The commutator direction $w$
is already zero in $H_1(L;\Z)$.
\end{proof}

\begin{remark}[Scope]
The numerical value $12$ is now an actual linking invariant in the
orbit five-sphere. The function $q$ alone cannot provide it: globally
$q$ is a meromorphic base coordinate with poles, so its argument
compares the zero and pole fibres. The globally defined function $h$
removes that ambiguity. The present result specifies one part of the
global attaching data for the cap of C.3.14; it does not evaluate its
Pontryagin framing term or its characteristic integer $J_\psi$.
\end{remark}
\begin{proposition}[C.3.17 (Proved): a marked circle extension over the cusp]
Let $V$ be a sufficiently small closed cusp body and
$\overline V=V/S^1_\delta$. Its smooth orbit structure can be
chosen compatibly with C.3.14 so that there is a diffeomorphism
\[
 \overline V\simeq E_{I_1}\times S^1,
\]
where $E_{I_1}$ is a closed neighbourhood of one nodal fibre in
an elliptic surface. More precisely, the actual positive
$\gamma$-translation circle action on the outer cusp collar
extends to a smooth circle action $\rho$ on $\overline V$.
The extension is free. It is not asserted to preserve the
fixed sphere $F\subset\overline V$ or the original interior
map to the $q$-disc.

Consequently the actual annular map $\psi$ of C.3.10 is the
endpoint of an ambient isotopy of $Q$ whose restriction to $F$
is the loop of embedded spheres $t\mapsto\rho_t(F)$.
After the marked deformation used in the proof, its trace is
the normalization immersion of a nodal elliptic curve times a
circle: it has one transverse double circle and two distinct
preimage circles. The regular cusp three-torus with period
basis $(\gamma,u,w)$ represents a primitive generator of
$H_3(\overline V;\Z)$; in particular it does not bound a
four-chain inside the cusp body.
\end{proposition}
\begin{proof}
We first remove the bounded period units by a marked smooth
deformation. On the cusp write, as in C.3.5,
\[
 Z(q)=sB_0+C(q),\qquad
 B_0=\begin{pmatrix}0&1\\-1&0\end{pmatrix},\qquad
 q=e^{2\pi i s}.
\]
For $0\leq a\leq1$ replace $C(q)$ by $(1-a)C(q)$, keeping
the same $A_2$ triangulation and the real periods $w,\delta$.
This is a deformation of the actual toric period quotient,
not a replacement inferred from its homology. The two period
automorphisms still have the same integral translations of the
fan and commute; only their nowhere-zero holomorphic unit
multipliers change. They depend smoothly on $a$, and even
holomorphically on a small complex neighbourhood of this
interval. On a sufficiently small punctured $q$-disc the
imaginary period matrices are invertible uniformly in $a$:
their leading part is $(\operatorname{Im}s)B_0$, while the
remaining matrices are uniformly bounded.

The fixed toric charts, with this extra parameter, therefore
give a smooth family of the same Mumford cusp quotients.
There is no new stabilizer at the central fibre: a nonzero
period translates every height-one ray by a nonzero lattice
vector, and cannot preserve any finite cone of the fan.
The usual finite fundamental collection of toric charts gives
proper discontinuity, uniformly for the compact parameter
interval. The central family has the constant compact
normalization $dP_6$ with its smoothly varying opposite-side
identifications. Thus it has a proper representative over a
uniform smaller $q$-disc. One may see the uniformity directly
by taking a relatively compact neighbourhood of this compact
central family: its boundary is disjoint from $q=0$, and after
shrinking $|q|$ no fibre meets that boundary. The nonzero
fibres are connected compact tori, so the selected representative
contains the whole nearby fibre. Denote these compact smooth
cusp bodies by $V_a$.

Their $S^1_\delta$ actions form a smooth semifree family, with
a smooth family of fixed two-spheres. Choose an equivariant
Ehresmann transport with the following two prescribed local
forms. On the outer collar keep $q$ and all real coefficients
in the marked period basis fixed. This is a smooth transport
there and carries the actual $\gamma$-translation action to
the corresponding action for every $a$. Near the fixed spheres,
choose family equivariant normal tubular coordinates. After
conjugating the negative-weight line the normal representation
is scalar multiplication on $\mathbb C^2$. Choose Hermitian
metrics and a unitary connection, and make the transport
linear and unitary in these normal coordinates. The collar
and fixed tubes are disjoint. A circle-invariant partition of
unity extends and patches the prescribed horizontal fields;
each has derivative one on $a$, so their patched field does
also. Compactness gives its flow for the entire interval.

This particular transport induces a smooth transport of orbit
spaces. On the free locus this follows from the principal-bundle
quotient. In the fixed tubes, unitary normal transport induces
the usual smooth linear rotation in the Hopf $\mathbb R^3$
coordinates. Thus it is smooth there as well. This explicit
choice avoids any assertion that every equivariant smooth map
has a smooth map in an arbitrarily preassigned quotient atlas.
The Hopf smooth structures and tubular choices are legitimate
ones for the orbit construction in C.3.14; the annular map is
supported in the free collar, where these structures coincide.
We obtain a marked orbit diffeomorphism
\[
 (\overline V,\text{boundary }\gamma\text{-action})
 \simeq(\overline V_1,\text{boundary }\gamma\text{-action}).
\]
The transport is not required to preserve $q$ in the interior.

At $a=1$ the period matrix is $sB_0$. The complex torus is
the product of two Tate elliptic curves, and the square toric
model is the fibre product of their nodal families. The chosen
$A_2$ triangulation is its chosen small resolution. Project
this resolution to the first Tate family $E_{I_1}$, and let
$S^1_\delta$ act on the second factor. The projection descends
to a proper smooth submersion
\[
 p:\overline V_1\longrightarrow E_{I_1}
\]
with oriented circle fibres. Here is the local check at the
only place where this assertion is not immediate. Write the
fibre product as $xy=uv$, with $q=xy=uv$, and take the small
resolution chart
\[
 u=xA,\qquad y=Av,
\]
whose coordinates are $(x,A,v)$. The weights of $(A,v)$ are
$(+1,-1)$ and $x$ has weight zero. Its Hopf quotient coordinates
are
\[
 \bigl(x,y=Av,d=(|A|^2-|v|^2)/2\bigr),
\]
and the descended projection is simply $(x,y,d)\mapsto(x,y)$.
The second chart, $x=uB$, $v=By$, gives the identical check
with $d'=(|u|^2-|B|^2)/2$. The local orbit fibre at $x=y=0$
consists of the two orbit rays joined by the quotient interval
of the exceptional projective line, and is a smooth real line.
Globally it closes to the second Tate orbit circle. At fixed
points away from this chart the model is likewise
$(p,A,v)\mapsto(p,Av)$, whose quotient is the same smooth
projection. We make the compatibility of these quotient charts
explicit. On overlaps near the same fixed branch the two weight
coordinates have the form
\[
 A_j=a_{ij}(p,z)A_i,\qquad v_j=b_{ij}(p,z)v_i,
 \qquad z=A_iv_i,
\]
where the units and the change of the tangential coordinate
$p$ are pulled back from the projection base. The two fixed
lifts over the node have disjoint fixed tubular neighbourhoods;
they require no overlap identification at the node itself.
Choose Hermitian metrics for the two weight line bundles on
these base charts, retaining the separated lifts of the fixed
normalization. If their positive local coefficients are
$h_{A,i},h_{v,i}$, their transition rules are
\[
 h_{A,j}|a_{ij}|^2=h_{A,i},\qquad
 h_{v,j}|b_{ij}|^2=h_{v,i}.
\]
The weighted Hopf coordinate
\[
 d_i=\tfrac12\bigl(h_{A,i}|A_i|^2-h_{v,i}|v_i|^2\bigr)
\]
therefore agrees with $d_j$ on overlaps. These are legitimate
Hopf coordinates: rescale the two normal variables by
$\sqrt{h_{A,i}}$ and $\sqrt{h_{v,i}}$. Their product is
$\sqrt{h_{A,i}h_{v,i}}\,z$, a smooth local change of the
complex coordinate $z$ with invertible derivative at $z=0$.
Thus $(p,z,d_i)$ are smooth quotient coordinates, their
overlap maps are smooth, and the projection remains
$(p,z,d_i)\mapsto(p,z)$. This specifies the Hopf smoothing
used for the transport above, without a general smoothness
claim for arbitrary equivariant quotient maps.

The section of the second Tate family represented by $z_2=1$
avoids its node. It therefore gives a section of the fibre
product which lifts through the small resolution, and then
gives a section of $p$. Hence this oriented circle bundle is
trivial. This proves the asserted product diffeomorphism by
an actual projection and section.

We can match the boundary action exactly, not just its homology
class. On the boundary of this circle bundle the positive
$\gamma$ action is the circle translation in the second Tate
orbit fibre. Its invariant vertical length form has fibre
integral one. Extend that length form to any smooth positive
vertical length form on the entire oriented circle bundle,
using a collar cutoff. Divide on each fibre by its integral;
this is a smooth positive function of the base because $p$
is a proper submersion. The dual vertical vector field now has
period exactly one, and agrees with the specified vector field
on the boundary collar. Its flow defines the desired free
circle action. Transport it back to $\overline V$ by the marked
orbit diffeomorphism.

Choose an invariant function $\chi$ equal to one in the inner
body and zero outside the outer collar, with the same collar
profile as in C.3.10. Such a function can be chosen because
the action agrees with the actual $q$-fibre translations there;
extend it constantly by one on the remaining interior. The
maps
\[
 H_t(z)=\rho_{t\chi(z)}(z),\qquad 0\leq t\leq1,
\]
are diffeomorphisms, since $\chi$ is invariant. They extend
by the identity to $Q$. Their endpoint is the actual annular
map: on the inner body it is $\rho_1=\id$, and on the collar
it is translation by $\chi\gamma$. On the fixed sphere
$\chi=1$, so this isotopy traces precisely $\rho_t(F)$.
It is not an isotopy relative to the fixed sphere.

In the zero-unit model the fixed sphere projects to the
normalization of the first nodal elliptic curve. Over a smooth
point of that curve it meets the orbit circle once, while the
two preimages of its node give two distinct points of the
orbit circle over the node. Its circle sweep is therefore the
normalization immersion of the nodal curve times the circle.
The local form $(x,y,d)$ with $q=xy$ gives its two transverse
branches at the double circle. Finally $E_{I_1}$ retracts onto
its nodal elliptic curve, whose fundamental two-cycle is the
degree-one specialization of a regular elliptic fibre. The
K\"unneth theorem now gives
\[
 H_3(E_{I_1}\times S^1;\Z)
 =H_2(E_{I_1};\Z)\otimes H_1(S^1;\Z)=\Z,
\]
generated by that regular fibre times the marked circle.
Boundary marking transports this three-torus to the actual
$(\gamma,u,w)$ torus, proving the last assertion.
\end{proof}

\begin{remark}[Scope]
The raw fractional period translations need not extend across
the central fibre. The action above modifies their interior
orbits. Its loop moves the fixed sphere, which is exactly the
feature retained by the annular relative mapping class. Thus
the construction does not trivialize that class or evaluate
$J_\psi$. In particular the nonbounding three-torus prevents
discarding its stable framing contribution by a chain lying
entirely in the cusp body.
\end{remark}
\begin{proposition}[C.3.18 (Proved): an actual bounding four-manifold and its frame defect]
There is a smooth genus-one Lefschetz fibration
\[
 \mathfrak g:Q\simeq S^5\longrightarrow S^3
\]
whose critical set is a circle mapping bijectively to the positive
$(3,4)$ knot
\[
 K=\{z_2^3-z_1^4=0\}\subset S^3.
\]
It agrees on the exterior and the outer cusp collar with the
actual torus quotient of C.3.13. In particular, the inverse image
of a nearby positive weighted circle $K'$ is the actual marked
three-torus with periods $(\gamma,u,w)$.

This three-torus bounds an embedded compact four-manifold
$M\subset Q$ carrying a positive genus-one Lefschetz fibration
over a Seifert surface of $K'$, with exactly twelve nodal critical
points. Its Euler number is $12$. Give $M$ its Lefschetz complex
orientation and its boundary the induced orientation, which in
the marked period convention is $(\gamma,w,u)$. Let
$\mathfrak f_{\mathrm{aff}}$ be the stable boundary frame consisting
of the three period directions and the two real $q$-normal
directions, with a fixed ordering consistent with the ambient
orientation. Relative to this actual affine frame,
\[
 \left\langle p_1(TM\oplus\R,\mathfrak f_{\mathrm{aff}}),
                   [M,\partial M]\right\rangle=-24.
\]
Equivalently, its transgression is $+24$ when the boundary
three-torus is oriented by the ordered basis $(\gamma,u,w)$.
No value of $\langle p_1(M_\psi)x,[M_\psi]\rangle$ or of
$J_\psi$ is inferred without the cap transgression calculation.
\end{proposition}
\begin{proof}
On the complement of the cusp body, the descended map of C.3.13
is a smooth two-torus bundle: on an ordinary chart it is the
quotient by the two residual real period directions $u,w$,
and the finite root charts give the same submersion after their
free finite quotients. The explicit finite gauges in C.3.13
ensure that this includes the two exceptional weighted circles
in the target.

Use the marked product of C.3.17 on the cusp body. In its
zero-unit model the projection is an elliptic $I_1$ neighbourhood
times a circle. Choose the product circle coordinate by the
second Tate identity section $z_2=1$ and the extending circle
action. On the boundary this coordinate is exactly the marked
character $\theta$: the section has value zero and the boundary
action is the actual $\gamma$ translation. The prescribed period
transport in C.3.17 preserves all these coordinates on the
outer collar. Consequently the map
\[
 (x,\theta)\longmapsto
 \Phi_0\bigl(q_E(x),\theta-\operatorname{Re}\kappa_0(q_E(x))\bigr)
\]
agrees there with the exterior map, where $q_E:E_{I_1}\to\Delta$
is the nodal elliptic map and $\Phi_0,\kappa_0$ are the exact
solid-torus coordinates of C.3.13. The two maps glue smoothly.
Near the critical circle, after the indicated smooth target
coordinate change, their local expression is
\[
 (x,y,\theta)\longmapsto(xy,\theta).
\]
Thus there is one positive nodal degeneration over each point
of $K$, and no other critical point. This constructs the stated
map from the actual marked collar. It need not equal the old
map $\overline f$ in the interior of the cusp body.

Choose $0<\epsilon\ll1$ and put
\[
 K'=\{(1-\epsilon)z_2^3-z_1^4=0\}\subset S^3,
 \qquad F_\epsilon=(1-\epsilon)z_2^3-z_1^4.
\]
Then $K'$ is the positive weighted circle $q=\epsilon$ in
the actual collar coordinates. The weighted circle field
$(z_1,z_2)\mapsto(e^{6\pi it}z_1,e^{8\pi it}z_2)$ satisfies
\[
 F_\epsilon(t\cdot z)=e^{24\pi it}F_\epsilon(z).
\]
Hence $\arg F_\epsilon$ is a submersion off $K'$. The closure
of one of its pages is a smooth compact oriented Seifert
surface $\Sigma'$ with boundary $K'$: the polynomial has a
simple complex normal zero along the link, which gives the
usual smooth half-plane collar of a page. Its fibres are
connected, since the normal meridian has argument degree one.
On the critical knot $K$,
\[
 F_\epsilon=-\epsilon z_2^3.
\]
Its argument increases twelve times along the positive weighted
parameter. Every page therefore meets $K$ transversely in
exactly twelve points, all with positive local intersection.

Set $M=\mathfrak g^{-1}(\Sigma')$. The displayed local nodal
model and transversality of $\Sigma'$ to the critical circle
show that $M$ is a smooth embedded compact four-manifold with
boundary $\mathfrak g^{-1}(K')$. The restriction to $M$ is a
positive genus-one Lefschetz fibration over $\Sigma'$ with
twelve critical points. A smooth torus has Euler number zero,
and an $I_1$ fibre has Euler number one. Stratification over
these twelve values therefore gives $\chi(M)=12$.

We next specify the boundary framing whose characteristic number
is computed. Choose a fibrewise complex structure on the regular
two-tori, compatible with their affine markings, and the usual
complex structure in the nodal charts. Such choices glue: on
the intervening collars the marked elliptic upper half-plane
is contractible. Choose a complex structure on $\Sigma'$ and
a compatible horizontal distribution. These give an almost
complex structure on $M$, equal to the standard nodal structure
near its critical points. Along $\partial M$, its vertical
complex tangent is framed by the translation vector dual to
the elliptic one-form with $w$-period one. Frame its horizontal
complex line by a lift of the outward normal of $\Sigma'$;
the associated real imaginary direction is the positive
boundary tangent. Denote this complex frame by $\mathfrak f_C$.

For completeness, the relative determinant class is pulled
back from the base. On a regular fibre its nonzero holomorphic
one-forms form a line $\lambda$, and evaluation identifies
the relative canonical line with the pullback of $\lambda$.
This remains true through each nodal chart: for $t=xy$ the
relative form is $dx/x=-dy/y$, and wedging with $dt$ gives
a nowhere-zero multiple of $dx\wedge dy$. Thus, as complex
line bundles,
\[
 \det(TM,J)\simeq\pi^*(T\Sigma'\otimes\lambda^{-1}).
\]
The determinant trivialization supplied by $\mathfrak f_C$
is the pullback of the corresponding base trivialization along
$\partial\Sigma'$. Therefore for some
$\beta\in H^2(\Sigma',\partial\Sigma';\Z)$,
\[
 c_1(TM,\mathfrak f_C)=\pi^*\beta,
 \qquad c_1(TM,\mathfrak f_C)^2=0.
\]
This does not require assigning a degree to $\lambda$ or a
holomorphic structure to the entire four-manifold. Its local
evaluation argument concerns the fibrewise elliptic structures
and the positive nodal charts.

The horizontal vector in $\mathfrak f_C$ points outwards on
the boundary. The relative top Chern class is therefore the
relative Euler class for an outward-pointing vector, and
Poincar\'e--Hopf gives
\[
 \langle c_2(TM,\mathfrak f_C),[M,\partial M]\rangle
 =\chi(M)=12.
\]
The relative characteristic-class identity now gives
\[
 \langle p_1(TM,\mathfrak f_C),[M,\partial M]\rangle
 =\langle c_1^2-2c_2,[M,\partial M]\rangle=-24.
\]

We finally compare this frame with the actual annular one.
On the target collar the exact formulas of C.3.13 give
\[
 F_\epsilon=(q-\epsilon)S(q)^{12}
              \exp(24\pi i\vartheta),
 \qquad\vartheta=\theta-\operatorname{Re}\kappa_0(q).
\]
Along $q=\epsilon$, the inward normal ray of the chosen page
therefore has $q$-angle equal to a constant minus $12\theta$.
Its complementary normal line in $S^3$ rotates by the same
angle. Add the pullback of that normal real line to $TM$,
identifying $TM\oplus\R$ with $TQ|_M$ after a horizontal
choice. On the boundary, the complex vertical translation
frame and the positive $\gamma$ tangent are the actual marked
period frame. The two remaining real vectors differ from
the fixed real and imaginary $q$ directions by a rotation
of winding $-12$ as $\gamma$ traverses the circle. Changes
of horizontal lifts add only contractible shear matrices.
Thus the stable frame comparison, up to fixed ordering, is
\[
 \operatorname{diag}\bigl(I_3,R_{-24\pi\theta}\bigr):
 T^3\longrightarrow SO(5).
\]
It factors through the $\gamma$ circle and is nullhomotopic,
since $\pi_1(SO(5))=\Z/2$ and twelve is even. Consequently
the relative first Pontryagin number is unchanged when
$\mathfrak f_C\oplus1$ is replaced by
$\mathfrak f_{\mathrm{aff}}$.

For the orientation convention, the elliptic complex orientation
is $(w,u)$, because its marked periods are $1,\tau$ with
$\operatorname{Im}\tau>0$. The outward-first boundary convention
for the Lefschetz orientation therefore orients the boundary
by $(\gamma,w,u)$. Reversing this to $(\gamma,u,w)$ reverses
the transgression sign, giving the equivalent $+24$ convention
stated above.
\end{proof}

\begin{remark}[Scope]
The four-manifold leaves the cusp body; its existence is
consistent with the local nonbounding assertion in C.3.17.
The value computed here is a relative Pontryagin defect of
the specified affine frame, not yet the Pontryagin pairing
of the closed cap in C.3.14. That latter conversion must
retain the actual gluing and its line-bundle transition.
\end{remark}
\begin{proposition}[C.3.19 (Proved): the cap Pontryagin pairing]
For the actual cap $M_\psi$ and meridian class $x$ of C.3.14,
\[
 \bigl|\langle p_1(M_\psi)x,[M_\psi]\rangle\bigr|=24.
\]
Thus the framing coordinate $\beta'$ in Wall's convention has
$|\beta'|=6$. The orientation sign and the cubic
$k_\psi=\langle x^3,[M_\psi]\rangle$ are not determined by
this proposition, and no numerical $J_\psi$ or monodromy class
is asserted.
\end{proposition}
\begin{proof}
We use the cohomology suspension of the actual universal class
$p_1$, denoted $s(p_1)\in H^3(SO;\Z)$, with the clutching
convention fixed once throughout. In particular we do not
replace it by an unspecified primitive generator of
$H^3(SO;\Z)$ or of $H^3(\operatorname{Spin};\Z)$.
All computations below may first be made in real cohomology;
they determine the integral characteristic number, and torsion
cannot contribute to its pairing.

Write $N=S^2\times D^4$ and $Y=\partial N=S^2\times S^3$,
so that the cap is $N\cup_{\widetilde\psi}(-N)$. Choose a
stable tangent frame on $N$. Its restriction supplies a stable
frame on the outer copy $P=S^2\times D^3\subset Y$. The
stable frame on $Q=S^5$, restricted to $P$, is homotopic to
this one: the difference is a map $P\to SO$, and
$P\simeq S^2$ while $\pi_2(SO)=0$. Thus the frame comparison
period obtained in C.3.18 is also the one computed from the
chosen frame on $N$.

Let $\mathcal A=L\times I\subset P$ be the annular support,
where
\[
 L=S^1_\gamma\times\operatorname{Nil}^3(u,w,q).
\]
Use a global affine nilmanifold frame on $\mathcal A$.
The nontrivial $q$-monodromy in $u,w$ is accounted for by
the invariant nilmanifold frame; on any fixed $q$-angle
three-torus it differs from the marked period frame only
by a constant triangular change. Translation in $\gamma$
preserves this global frame. Let
$A:\mathcal A\to SO$ be the stable change-of-frame map
from this frame to the restricted global frame.

In the affine frame, the derivative of the annular map
$(r,p)\mapsto(r,\rho_{\chi(r)}p)$ is the identity plus
the single shear $\chi'(r)\,dr\otimes\partial_\gamma$.
Multiplying the shear by a parameter from one to zero
contracts it to the identity through invertible matrices,
fixed on both end collars. Therefore the actual derivative
clutching map is homotopic, relative to the complement of
$\mathcal A$, to
\[
 G(r,p)=A(r,\rho_{\chi(r)}p)^{-1}A(r,p).
\]
At both ends $G=1$, since the translation is respectively
zero or a full period. Replacing $A(r,p)$ by its restriction
to one middle slice gives a homotopy of $G$ relative to
these ends. We may consequently write $A=A(p)$.

The Mayer--Vietoris clutching formula gives
\[
 p_1(M_\psi)=\delta\bigl(G^*s(p_1)\bigr),
\]
where $G$ is extended by the identity on the rest of $Y$.
With the boundary orientation associated to this connecting
map, its cup-product pairing is exactly
\[
 \langle p_1(M_\psi)x,[M_\psi]\rangle
 =\langle G^*s(p_1)\smile x|_Y,[Y]\rangle.
\]
Equivalently this is the relative pairing on
$(\mathcal A,\partial\mathcal A)$. Changing a global
clutching orientation reverses both associated signs;
we will retain the absolute value in the numerical conclusion.

We calculate this relative class explicitly. Define
\[
 \Gamma:S^1\times L\longrightarrow SO,
 \qquad\Gamma(t,p)=A(\rho_t p)^{-1}A(p).
\]
It is the identity on $\{0\}\times L$. The suspension
$s(p_1)$ is primitive under multiplication. Slanting its
pullback by the oriented $t$-circle therefore gives
\[
 \int_{S^1}\Gamma^*s(p_1)
 =-\int_{S^1}\rho^*\bigl(A^*s(p_1)\bigr)
 \quad\hbox{in }H^2(L;\R).
\]
There is no additional relative lift ambiguity here:
restriction $H^2(S^1\times L)\to H^2(\{0\}\times L)$
is surjective. Thus $H^3(S^1\times L,\{0\}\times L)$
is exactly the degree-two slant summand. Pulling back by
$t=\chi(r)$ identifies the relative annular class with
the suspension of the right-hand side. This is the
relative trace calculation, including the full-period
endpoint; an absolute equality on $\mathcal A$ alone
would not suffice.

Put $g=\theta_\gamma$ and $a=\arg q/(2\pi)$. The meridian
Euler class has, up to the already fixed overall circle
orientation, the representative
\[
 e=x|_L=dg\wedge da.
\]
Indeed the positive $q$ period has
$A_0\gamma=\gamma-\delta$ and $A_0u=u+w$.
Using coefficient coordinates on the covering mapping
torus, its deck identification can be written
\[
 (g,u,w,d,a)\longmapsto(g,u,w-u,d+g,a+1),
\]
where $d$ is the $\delta$ coordinate. The normalized
connection $\mathrm d d-a\,\mathrm d g$ is invariant and has curvature
$dg\wedge da$. This also shows that no additional
component of the Euler class has been suppressed.

Decompose the real cohomology class
$\alpha=A^*s(p_1)$ using $L=S^1_g\times\operatorname{Nil}^3$:
\[
 \alpha=dg\smile\beta+\alpha_0,
 \qquad\beta\in H^2(\operatorname{Nil}^3;\R),\quad
 \alpha_0\in H^3(\operatorname{Nil}^3;\R).
\]
Translation in $g$ has trace $\beta$. Therefore the preceding
relative calculation and the cup pairing reduce, with a
fixed orientation sign, to
\[
 \pm\int_L dg\wedge da\wedge\beta
 =\pm\int_{T^2_{u,w}}\beta
 =\pm\int_{T^3_{\gamma,u,w}}A^*s(p_1).
\]
The first equality on the right is integration over the
mapping-torus base, with its fibre orientation. Components
of $\beta$ containing $da$, and the component $\alpha_0$,
give zero. This proves that there are no unaccounted
cross terms in the cap pairing.

Finally, the last integral is the relative stable-frame
Pontryagin defect of C.3.18, with at most the inverse
change-of-frame and orientation signs. Its absolute value
is $24$. Hence the cap pairing has absolute value $24$.
Wall's exact formula $\langle p_1x,[M_\psi]\rangle=4\beta'$
from C.3.14 gives $|\beta'|=6$, as asserted.
\end{proof}

\begin{remark}[Scope]
The argument uses both the global frame defect furnished by
the bounding four-manifold and the actual Euler class of the
$\delta$ bundle on the annulus. A local unipotent derivative
alone would give neither. The cubic pairing and its orientation
relative to this Pontryagin pairing remain necessary to
evaluate Wall's unframed knot coordinate.
\end{remark}
\begin{proposition}[C.3.20 (Proved): the cubic term as an actual relative square]
The four-manifold in C.3.18 can be chosen disjoint from the fixed
sphere. Write it as $M\subset P$, retain its Lefschetz orientation,
and let $\mathcal L$ be the actual principal $\delta$ circle bundle
on $P$, regarded as a complex line bundle. Its restriction to
$\partial M=T_{\gamma uw}$ is trivial. Put
\[
 s_\delta=\langle c_1(\mathcal L|_M,\mathrm{triv})^2,
                         [M,\partial M]\rangle.
\]
This integer is independent of the choice of boundary trivialization.
For the cap of C.3.14, there is one overall sign $\epsilon_{\rm cap}$,
coming from its comparison with the oriented cut-and-reglue trace,
such that
\[
 k_\psi=\epsilon_{\rm cap}\,3s_\delta.
\]
The same relative square, with the compatible Hopf-meridian
orientation, is the linking coefficient of the marked
$\delta=0$ lift of $T_{\gamma uw}$ in $X\setminus F\simeq S^3$.
No numerical value of $s_\delta$ is assumed in this statement.
\end{proposition}
\begin{proof}
First we arrange the claimed disjointness. In the marked product
of C.3.17, the section of the second Tate family avoids $F$.
Consequently the product circle coordinate of the compact sphere
$F$ has image in a compact proper arc of $S^1\setminus\{0\}$.
An orientation-preserving reparametrization of this circle can
compress that arc to length less than $1/12$. Perform this
reparametrization over a neighbourhood of the central first
elliptic fibre, interpolate through circle diffeomorphisms to the
identity, and leave the outer collar fixed. This changes only the
interior circle coordinate and the extending action, not the
actual annular endpoint $\psi$ or its boundary marking.

Rebuild the map of C.3.18 with this coordinate. The twelve
intersection parameters of one Milnor page with the critical
knot are equally spaced in the product circle. Choose the phase
of the page so that all twelve avoid the compressed arc.
Since $F$ lies over the central elliptic fibre, the inverse image
$M$ of this page is now disjoint from $F$. It has the same
boundary, twelve positive nodes, orientation and affine frame
as in C.3.18. Shrink the deleted fixed tube so that $M\subset P$.

Let $g,u,w$ denote the marked circle coordinates on a positive
real regular cusp fibre, let $a$ be the positive argument of
$q$, and let $r=|q|$. Coefficient coordinates have deck rule
\[
 (g,u,w,d,a)\longmapsto(g,u,w-u,d+g,a+1).
\]
Thus the invariant normalized circle connection is
$\alpha=dd-a\,dg$, with Euler curvature
\[
 E=dg\wedge da.
\]
Choose a global connection on $\mathcal L$ agreeing with this
one on the annular collar. This is possible because the displayed
connection is on the actual restricted bundle. Also choose its
curvature near $\partial P$ to be the pullback of the Hopf
curvature on the normal two-sphere. Then $E^2=0$ near that
boundary. Since $H^4(P)=H^3(P)=0$, there is a three-form $B$
with $dB=E^2$, unique up to an exact form. Because
$H^3(\partial P)=0$, it can moreover be chosen zero on a
boundary collar.

The line bundle on the positive real torus is trivialized by
the marked $d=0$ section. On that torus the displayed connection
is zero in this section. Hence relative Chern--Weil theory and
Stokes' theorem give, with $\partial M=-T_{\gamma uw}$,
\[
 s_\delta=\int_M E^2
          =-\int_{T_{\gamma uw}}B.
\]
Changing the trivialization changes the relative Chern class by
the connecting image of $H^1(\partial M;\Z)$. Such a class is
in the radical of the relative intersection pairing: its
pairing with a relative class is the boundary pairing with the
restriction of that class's absolute image, which is zero here.
Products of two connecting classes are zero as well. Thus the
square is independent of the boundary gauge.

We calculate the cut-and-reglue variation of $x^3$. Double $P$
to $B_5=DP=\partial N$, where $N=S^2\times D^4$, and extend
$\psi$ and all variations by the identity on the second half.
Here is an explicit seven-dimensional bordism used for this
comparison. Form the $B_5$-bundle over an annulus with monodromy
$\widetilde\psi$. Divide its outer boundary circle into four
successive arcs. Attach two copies of $N\times[0,1]$ along the
$B_5\times[0,1]$ faces over opposite arcs, reversing the interval
parametrization on the second attachment. Choose the
trivializations over these two arcs so that transport over one
remaining arc is the identity and transport over the other is
$\widetilde\psi$. After rounding corners, those remaining arcs
join the endpoint copies of $N$ into, respectively, the identity
double and the $\widetilde\psi$-twisted double. The inner boundary
is the mapping torus. Thus the three closed boundary components
are the two actual caps and the mapping torus, with the two caps
oppositely oriented. Their comparison has a single orientation
sign fixed by these conventions, independent of the map.

The meridian line extends over this actual bordism: use its
specified lift of $\widetilde\psi$ on the annular bundle and
the pullback of its line on $N$ on each attached $N\times[0,1]$.
The attaching identifications are precisely the chosen line
identifications, so there is no additional clutching datum.
The tangent bundle is the tangent bundle of this smooth
seven-manifold; on each boundary it restricts stably to that
boundary's tangent bundle. Consequently the same bordism
compares $x^3$ and $p_1x$, with the same orientation sign.
The identity cap is $S^2\times S^4$, and the identity mapping
torus has a line pulled back from $B_5$; both characteristic
pairings vanish on both identity manifolds. Equivalently one
may subtract the identity mapping torus, which bounds
$B_5\times D^2$ with its pulled-back line. Thus the desired
cubic is the mapping-torus variation, up to this one common
boundary-orientation sign.

On the annulus $\psi:g\mapsto g+\chi(r)$, where $\chi$ decreases
from one to zero. Consequently
\[
 \Delta=\psi^*E-E=d\chi\wedge da.
\]
There is a global one-form $\eta$ with $d\eta=\Delta$ whose
annular value is
\[
 \eta=-(1-\chi)\,da-12\,dg.
\]
Indeed the outer value $-(da+12\,dg)$ extends as a closed form
over the exterior by C.3.16, while the inner value $-12\,dg$
extends over the product cusp by C.3.17. Closed representatives
can be made to agree exactly on collars by subtracting exact
forms. An exact correction makes $\eta$ zero near $\partial P$,
since this boundary has $H^1=0$. This correction does not affect
any integral below, as $B$ vanishes there.

We use the normalized characteristic-form transgression of
\citet[Section 3, Propositions 3.2 and 3.8]{ChernSimons1974E}.
The interpolating curvature on the mapping cylinder is
$E+t\Delta+dt\wedge\eta$. Its endpoint connections can be
matched to the chosen clutching lift by an exact gauge correction,
because $H^1(B_5;\R)=0$. The transgression computation gives
\[
 \int(E+t\Delta+dt\wedge\eta)^3
 =\int_P\eta\wedge(3E^2+3E\Delta+\Delta^2)
 =3\int_P\eta\wedge E^2.
\]
Here $E\Delta=\Delta^2=0$ on the support of $\Delta$. Since
$d(\eta\wedge B)=\Delta\wedge B-\eta\wedge E^2$,
the last expression is $3\int_P\Delta\wedge B$.
Use the orbit orientation
$dg\wedge du\wedge dw\wedge dr\wedge da$, obtained by
taking the positive $\delta$ direction first in the zero-unit
complex model. Integration over the annulus, with $\chi$ going
from one to zero, now gives
\[
 \int_P d\chi\wedge da\wedge B
   =-\int_{T_{\gamma uw}}B=s_\delta.
\]
This proves the asserted factor of three and fixes it without
using a primitive-generator normalization in $H^3(SO)$.

Finally on the total circle bundle over $P$ the form
\[
 \Omega=\alpha\wedge\pi^*E-\pi^*B
\]
is closed. Its integral on the normal Hopf three-sphere is one:
the first term is the normalized Hopf integral and the second
vanishes in the chosen fixed collar. It therefore represents
the integral generator of $H^3(X\setminus F;\Z)$. On the marked
lift of the positive real $T_{\gamma uw}$, the connection and
Euler curvature restrict to zero, and its integral is
$-\int_T B=s_\delta$. This is precisely the claimed linking
coefficient. In particular the square is an invariant of the
actual inclusion, not merely of the abstract genus-one surface.
\end{proof}

\begin{remark}[Scope]
To evaluate this number by a height calculation one must identify
the actual line bundle $\mathcal L|_M$, including its mixed
base--fibre class, with the Poincar\'e line attached to the
original normal-function section. The numerical square of an
auxiliary split or torsor surface does not make that
identification. The signed cancellation with the Pontryagin
term likewise uses the same cut-and-reglue orientation in both
calculations.
\end{remark}
\begin{proposition}[C.3.21 (Proved): the actual delta-line square]
For the actual four-manifold and line bundle in C.3.20,
\[
 s_\delta=-2.
\]
The comparison is with the actual relative Picard section of
the circle bundle; it does not assume that the compact-core
model of Part B fills the original six-sphere.
\end{proposition}
\begin{proof}
We first identify the base precisely. On the Milnor page used in
C.3.18, choose its phase so that
$F_\epsilon=(1-\epsilon)z_2^3-z_1^4$ is positive real.
Positive weighted rescaling identifies its interior with
\[
 C^\circ:\quad b^3=a^4+1,
 \qquad
 a=\frac{z_1}{F_\epsilon^{1/4}},\qquad
 b=\frac{(1-\epsilon)^{1/3}z_2}{F_\epsilon^{1/3}}.
\]
Conversely, every point of this affine curve has a unique
positive weighted rescaling to the unit three-sphere. A different
page phase changes this by constant complex factors and gives
the same marked cover. The compactification $C$ has genus three
and one point at infinity. Its actual map to the original
$t$-line is
\[
 t_\epsilon=\frac{a^4+1}{1+\epsilon a^4}.
\]
It is ramified with orders three and four over $0$ and $1$,
and with order twelve at the added point over $t=1/\epsilon$.
The twelve points over $t=\infty$ satisfy
$a^4=-1/\epsilon$ and are unramified. Thus this is the
actual page cover; the cover ramified twelvefold over the
original cusp is its $\epsilon=0$ comparator, not the same
marked map when $\epsilon>0$.

Keep the product modification inside $|q|<r_0<\epsilon$.
The page equation there has twelve disjoint graph discs
$\theta=\theta_j(q)$, since its derivative in $\theta$ is
twelve. Consequently $M$ consists of twelve nodal elliptic
neighbourhoods attached to an entirely untouched raw torus
quotient. Denote the latter part by $M^\circ$ and its base by
$C_0$; it is the page with these twelve discs removed. It has
nonempty boundary and retracts to a graph.

We identify the actual circle bundle on this unmodified part.
On an ordinary marked chart let $g$ be a real lift of its
$\gamma$ coefficient and set
\[
 \xi=\zeta_1-6\mu g,\qquad
 \eta=\zeta_2-\beta g.
\]
Then $\xi=\tau u+w$ and $\eta=\mu u+d$. Exponentiating
$\eta$ gives the actual degree-zero line on the first elliptic
torus with multipliers
\[
 (\xi,z)\sim(\xi+\tau,e^{2\pi i\mu}z),\qquad
 (\xi,z)\sim(\xi+1,z).
\]
These are the Poincar\'e multipliers of the original normal
function $P$, up to the single dualization convention for the
associated circle line. The identification is an identification
of the varying fibrewise Picard section, not just its fibre
degree. Under a change of marked elliptic periods the maps are
fibrewise affine holomorphic and the line maps are its
semiabelian lifts. The translations may depend smoothly on
$g$; no global holomorphic-torsor assertion is needed.

Over the graph base $C_0$, the affine elliptic torsor has a
smooth section. Indeed its obstruction lies in the twisted
group $H^2(C_0,\Lambda)$, which is zero; equivalently use the
sequence from the lattice local system to its real vector
bundle and then to the torus bundle, whose vector-bundle
section sheaf is fine. Identify the torsor with its Jacobian
using such a section. Translation of a degree-zero Poincar\'e
line changes it by a line pulled back from the base. Two
fibrewise holomorphic lines with the same varying Picard
section likewise differ by a base line: their quotient has a
one-dimensional space of fibrewise holomorphic sections and
evaluation trivializes it fibrewise. Since $H^2(C_0;\Z)=0$,
this base line has zero topological Chern class. The actual
class on $M^\circ$ is therefore the Poincar\'e class of
$P\circ t_\epsilon$, up to sign, including its mixed
base--fibre component.

There is no hidden correction supported at the twelve modified
nodes. By excision and Poincar\'e--Lefschetz duality, each
local summand of
\[
 H^2(M,M^\circ;\Z)
 \simeq\bigoplus_{j=1}^{12}H^2(E_{I_1},\partial E_{I_1};\Z)
\]
is $\Z$, generated by the Thom class of a regular elliptic
fibre. The regular fibre is primitive in $H_2(E_{I_1};\Z)$,
as its degree-one specialization generates the nodal fibre's
two-cycle. Its image in $H^2(M;\Z)$ is the pullback of the
point class of the base, and is zero because that base has
boundary. The restriction $H^2(M;\Z)\to H^2(M^\circ;\Z)$
is consequently injective. This addresses actual local
clutching classes, rather than only their holonomy.

Trim a smaller boundary collar, which does not change the
relative square because its line is trivial, and identify the
remaining annular end of the page with the punctured disc at
the added point of $C$. Fill it by its regular marked
$T^2\times D^2$. The boundary affine torus bundle is
fibrewise complex-trivialized by a path of translations.
After matching the Picard section $P$, a remaining line
identification is a scalar base-line gauge. Thus one can choose
the product-cap Poincar\'e extension whose section is
$P\circ t_\epsilon$ near the added point; the displayed
multipliers are regular there. The line is topologically
trivial on the cap, and its square contribution is zero.
Boundary gauge changes cannot change the relative square, by
C.3.20. Call the resulting
closed smooth genus-one Lefschetz surface $\widehat M$ and its
line $\widehat{\mathcal L}$. Its square is the same relative
square $s_\delta$; the cap contributes zero.

We spell out why the square can now be compared with the
Jacobian without a projectivity or multisection assumption on
$\widehat M$. Let $j:C^*=C\setminus\{12\text{ nodal values}\}
\hookrightarrow C$ and let $\Lambda=R^1\pi_*\Z|_{C^*}$.
For a nodal elliptic fibre, degree-one specialization identifies
$R^2\pi_*\Z$ with the constant sheaf $\Z$, while degree one
is the invariant subgroup of the unipotent monodromy. Thus
\[
 R^1\pi_*\Z=j_*\Lambda,\qquad R^2\pi_*\Z=\Z.
\]
These sheaves and their fibre intersection pairing agree for
the actual surface and its Jacobian. The degree-zero class
$c_1(\widehat{\mathcal L})$ lies in the first Leray filtration
step. Its component in
$H^1(C,j_*\Lambda)$ is the connecting class of the varying
Picard section $\pm P$. This holds on $C^*$ by the preceding
actual multiplier calculation. It holds on all of $C$ because
\[
 H^1(C,j_*\Lambda)\longrightarrow H^1(C^*,\Lambda)
\]
is injective, by the first terms of the Leray sequence for $j$.
Local nodal choices can therefore add only a base-filtration
class; the excision argument gives the same conclusion directly.

Multiplicativity of Leray now determines the square. In degree
four its only nonzero graded piece is
$H^2(C,R^2\pi_*\Z)=\Z$. The square of a class in filtration
one is the cup square of its $H^1(C,j_*\Lambda)$ component,
using the fibre intersection pairing. Adding filtration-two
base classes contributes only filtration three or four, both
zero on this base. Hence the actual square is exactly the
square of the Poincar\'e class on the Jacobian; dualizing the
line does not change it.

Finally let $J_\epsilon\to C$ be the relatively minimal
Jacobian obtained by pulling back the original elliptic surface
along $t_\epsilon$. It has the explicit affine model
\[
 Y^2=X^3-3\lambda bX+2\lambda^2a^2,
 \qquad P=(0,-\sqrt2\lambda a),\qquad
 \lambda=1-\epsilon.
\]
Writing $h=1+\epsilon a^4$, the substitution from the original
Weierstrass coordinates is
$x=(ab^2/h)^2X$, $y=(ab^2/h)^3Y$.
The discriminant is
\[
 \Delta=1728\lambda^3(1+\epsilon a^4).
\]
It has twelve simple zeroes on $C^\circ$, and $c_4$ is a unit
at all of them. They are precisely the twelve $I_1$ fibres.
At the unique point at infinity, $a,b$ have pole orders
three and four. Scaling $X,Y$ by a local parameter to orders
two and three gives a regular Weierstrass model whose
discriminant is a unit when $\epsilon\ne0$; the point at
infinity is therefore a good fibre. This also shows that the
fundamental line has degree one, so
$\chi(\mathcal O_{J_\epsilon})=1$.

The displayed section $P$ is finite on every affine chart and
also in this integral chart at infinity, hence it is disjoint
from the zero section. At a simple discriminant zero it is a
smooth point of the nodal fibre: a node has $Y=0$, whereas
$a\ne0$ there and $Y(P)=-\sqrt2\lambda a\ne0$.
All singular fibres are irreducible, so there is no component
correction. These explicit identities are independently checked
in \texttt{ab\_push\_actual\_line\_comparison.py}.

The divisor $P-O$ has fibre degree zero and represents the
required relative Picard section, up to a base twist. Both
sections have square $-\chi=-1$ and $P\cdot O=0$. Therefore
\[
 (P-O)^2=-1-1-2(P\cdot O)=-2.
\]
Base twists do not change this square. The actual Leray
comparison above proves $s_\delta=-2$, as asserted.
\end{proof}
\begin{lemma}[C.3.22 (Proved): one common signed cap convention]
Retain the relative Pontryagin number
$p_{\mathrm{rel}}=-24$ of C.3.18 and the actual relative
Euler square $s_\delta$ of C.3.20. With the single
cut-and-reglue orientation sign $\epsilon_{\rm cap}$ used
in C.3.20, the two cap pairings satisfy
\[
 \langle p_1(M_\psi)x,[M_\psi]\rangle
       =\epsilon_{\rm cap}\,p_{\mathrm{rel}},
 \qquad
 k_\psi=\epsilon_{\rm cap}\,3s_\delta.
\]
In particular
\[
 J_\psi=\epsilon_{\rm cap}
              \frac{p_{\mathrm{rel}}-12s_\delta}{24}.
\]
This lemma does not assign a value to $s_\delta$.
\end{lemma}
\begin{proof}
We fix the normalization by connections, so that a boundary
frame and a coordinate clutch are not silently identified.
For a real orthogonal connection with curvature $F$, use
\[
 p_1(F)=-\frac{\operatorname{tr}(F\wedge F)}{8\pi^2},
 \qquad
 \operatorname{CS}_{p_1}(\omega)
 =-\frac{\operatorname{tr}(\omega\wedge d\omega
                          +\tfrac23\omega^3)}{8\pi^2}.
\]
For $K$ a change of coordinates,
\[
 \operatorname{CS}_{p_1}(K^{-1}dK)
 =\frac{\operatorname{tr}(K^{-1}dK)^3}{24\pi^2}
 =:\tau_{p_1}(K).
\]
This is the suspension used for a mapping torus oriented
with its interval parameter first. Indeed interpolating
the connection from zero to $K^{-1}dK$ gives curvature
$dt\wedge\omega+(t^2-t)\omega^2$, whose first Pontryagin
form integrates in $t$ to the displayed three-form.

In contrast, let a boundary frame be
$\mathfrak f=\mathfrak t B$ relative to a global stable
frame $\mathfrak t$. The connection for which
$\mathfrak f$ is parallel has matrix $-dB B^{-1}$ in
the $\mathfrak t$ coordinates. Therefore
\[
 p_1(E,\mathfrak f)[M,\partial M]
 =-\int_{\partial M}\tau_{p_1}(B).
\]
Equivalently, capping the boundary frame by a standard
frame uses the inverse coordinate clutch $B^{-1}$.
This minus sign is forced by the connection computation;
it is not an orientation chosen to obtain a cancellation.

Now use exactly the change-of-frame convention of C.3.19:
\[
 \mathfrak t=\mathfrak f_{\mathrm{aff}} A,
 \qquad
 G=A(\psi p)^{-1}A(p)
\]
after contracting the affine derivative shear. Set
\[
 a_0=\int_{T_{\gamma,u,w}}\tau_{p_1}(A).
\]
Here $B=A^{-1}$ in the boundary-frame formula. Since
the boundary orientation of $M$ is
$\partial M=-T_{\gamma,u,w}$, that formula gives
\[
 p_{\mathrm{rel}}
 =-\int_{\partial M}\tau_{p_1}(A^{-1})
 =-a_0.
\]

For the actual annulus use the orbit orientation
$dg\wedge du\wedge dw\wedge dr\wedge da$ and the
Euler curvature $E=dg\wedge da$, with $\chi$ decreasing
from one to zero as $r$ increases. The relative trace
calculation of C.3.19 has representative
\[
 \tau_{p_1}(G)
 =-d\chi\wedge\beta
 \quad\hbox{in relative cohomology},
 \qquad
 \int_{T^2_{u,w}}\beta=a_0.
\]
Consequently
\[
 \int_P\tau_{p_1}(G)\wedge E=-a_0=p_{\mathrm{rel}}.
\]
For example, in the only contributing term the form is
$-a_0\,d\chi\wedge du\wedge dw\wedge dg\wedge da$;
relative to the displayed orbit orientation its coefficient
is $a_0\chi'(r)$, and the integral of $\chi'$ is $-1$.

This is also the Pontryagin--Euler pairing of the actual
mapping torus in C.3.20. Its line curvature restricts to
$E+t\Delta$ on the $t$-slices. The possible extra term
has integral zero, because
\[
 \int_{DP}\tau_{p_1}(G)\wedge\Delta
 =\int_{DP}\tau_{p_1}(G)\wedge d\eta=0.
\]
Thus no line-clutch correction changes the preceding
pairing. The same cut-and-reglue bordism of C.3.20 is a
smooth bordism of tangent bundles and of the specified
line bundle. Both characteristic polynomials $p_1x$
and $x^3$ therefore acquire the same boundary-orientation
sign $\epsilon_{\rm cap}$. Their identity-cap and
identity-mapping-torus values vanish. Combining the
computed mapping-torus pairing with the cubic calculation
of C.3.20 gives the two asserted signed equations and
then their stated combination.
\end{proof}
\begin{proposition}[C.3.23 (Proved): an additive integer before reduction modulo 28]
Fix the linear semifree circle action of C.3.14, with its
circle direction fixed. On the group of orientation-preserving
equivariant diffeomorphisms of $S^6$ there is an integer
\[
 \mathfrak q:\pi_0\operatorname{Diff}^+_{S^1}(S^6)
                    \longrightarrow\Z
\]
defined by the unframed Haefliger coordinate of the orbit
three-knot of equivariant clutching. It is a group homomorphism.
For a representative normalized to be the identity on a fixed
tubular neighbourhood, it is Wall's unframed knot coordinate,
so $\mathfrak q=-J$ in the convention of C.3.14. Its image in
the ordinary group $\Theta_7$ is the fixed signed reduction
modulo $28$ described there. In particular, if an equivariant
representative has a sixth power equivariantly isotopic to a
representative with $J=0$ as an integer, its ordinary mapping
class is zero.
\end{proposition}
\begin{proof}
We first make the normalization and its independence explicit.
The fixed set is $F=S^2$, and its equivariant normal bundle is
the trivial scalar complex rank-two bundle after conjugating
the negative weight, by C.3.14. An equivariant real derivative
on that normal representation is complex linear and has
positive real determinant. An orientation-preserving ambient
diffeomorphism therefore restricts to an orientation-preserving
diffeomorphism of $F$.

By \citet[Theorem A, page 621]{Smale1959E}, this restriction is
isotopic to the identity. Lift such an isotopy using the product
normal bundle and an invariant tubular cutoff, extending its
velocity field equivariantly. After this adjustment the map
fixes $F$ pointwise. Its normal derivative is a map
$F\to GL_2(\C)$. Polar decomposition and $\pi_2(U(2))=0$
give a homotopy to the identity, again realized in the fixed
tube by an equivariant isotopy. Finally the usual rescaling
of a tubular germ to its derivative, followed by an invariant
cutoff, makes the whole map the identity on a smaller fixed
tube. Equivariance rules out a derivative from the nontrivial
normal representation to the fixed tangent representation;
the rescaled germ is smooth at the zero parameter. This
normalization uses equivariant isotopies throughout.

Clutch two standard linear equivariant seven-discs by this
diffeomorphism. Their fixed three-discs clutch to $S^3$.
Their orbit six-discs clutch to a homotopy six-sphere, which
is standard because $\Theta_6=0$. The fixed image is thus a
smooth oriented three-knot in $S^6$. Its Haefliger integer is
the proposed $\mathfrak q$, using the connected-sum group of
oriented three-knots in the six-sphere
\citep[Sections 1 and 5.16]{Haefliger1966D}.

This integer does not depend on the chosen normalization.
An equivariant isotopy of boundary clutching maps extends
over a collar to a diffeomorphism of the clutched equivariant
pairs. The resulting orbit pairs are diffeomorphic. One can
choose compatible unitary normal tubes along the entire
isotopy, so that their Hopf orbit charts form a smooth family;
this does not invoke functorial smoothness for arbitrary
equivariant maps in a fixed raw quotient atlas.
There is also no ambiguity from identifying the orbit homotopy
sphere with $S^6$. An orientation-preserving diffeomorphism of
$S^6$ is isotopic to one supported in a ball disjoint from a
given compact three-knot: make its germ the identity near a
point, use stereographic coordinates on the complement, and
conjugate its compact support into a small ball by translations
and positive dilations, all isotopic to the identity. Such a
representative fixes that knot. Consequently the ambient
identification cannot change its knot integer.

We prove additivity geometrically. Normalize two equivariant
diffeomorphisms to be the identity near $F$. In particular each
is the identity near a chosen fixed point. Stereographic
projection from that point identifies the complement with
$\R^2\oplus\C^2$, with the circle acting scalarly on the
second summand. The diffeomorphism has compact support there.
Positive dilations and translations along the fixed $\R^2$
commute with the circle and are equivariantly isotopic to the
identity on the one-point compactification. By conjugation,
place the two supports inside disjoint small invariant balls
centered at distinct fixed points. The resulting maps remain
in their original equivariant components.

For disjoint supports, the composition clutch is the
equivariant connected sum of the two separate clutches.
This follows directly by cutting the equatorial sphere between
the two invariant balls and inserting a product collar; the
map on the separating collar is the identity. The fixed
three-discs undergo the same boundary connected-sum operation.
On the orbit pairs this is the ordinary connected sum of
the two oriented three-knots. The defining group operation
of their Haefliger integers is connected sum. Hence
\[
 \mathfrak q(\phi_1\phi_2)
 =\mathfrak q(\phi_1)+\mathfrak q(\phi_2).
\]
This proves additivity without passing first to $\Theta_7$.

For a fixed-tube-relative representative, the surgery cap and
Wall calculation of C.3.14 apply verbatim: the inherited product
normal framing gives $\mathfrak q=-J$. Altering this framing
changes only Wall's framing coordinate, whose equivariant lift
extends over the seven-disc by the explicit quaternionic
formula in C.3.14. Finally the orbit-knot generator comparison
of \citet[Example 7.14 and Proposition 7.15]{Davis1982D}
identifies the underlying ordinary sphere class with the fixed
signed reduction of $\mathfrak q$ modulo $28$.

If $\phi^6$ is equivariantly isotopic to a representative with
integer $J=0$, then $6\mathfrak q(\phi)=0$ in $\Z$.
Thus $\mathfrak q(\phi)=0$, and its underlying ordinary
mapping class is zero. This last implication uses a torsion-free
integer before reduction; it is not inferred from the weaker
equation $6[\phi]=0$ in $\Z/28$.
\end{proof}
\begin{theorem}[C.3.24 (Proved): the actual smooth monodromy is trivial]
For the actual retained $S^6$ family and a positive generator
of its punctured $w$-parameter disc, the oriented smooth
monodromy satisfies
\[
 [\Phi]=0\quad\hbox{in }\pi_0\operatorname{Diff}^+(S^6)
                       \simeq\Theta_7\simeq\Z/28.
\]
Its mapping torus is consequently oriented-diffeomorphic to
$S^1\times S^6$. Every finite cyclic base change has trivial
smooth monodromy as well. By C.4.1's marked family comparison,
the same conclusion applies to the admissible sphere twist
data covered by that comparison.
\end{theorem}
\begin{proof}
Let $\widehat\psi$ be the actual smooth equivariant lift of
the annular orbit map in C.3.10.
The actual cap calculation is an integer calculation.
C.3.18 gives $p_{\mathrm{rel}}=-24$ and C.3.21 gives
$s_\delta=-2$. The common signed comparison in C.3.22 yields
\[
 J_\psi=\epsilon_{\rm cap}
          \frac{-24-12(-2)}{24}=0\quad\hbox{in }\Z.
\]
By C.3.14 and C.3.23, the unframed orbit-knot integer of the
actual annular lift is therefore zero. Neither the overall cap
orientation sign nor the Wall--Davis generator sign can change
this conclusion.

We verify that this is the sixth power at the equivariant
integer level. The family carries the fixed positive
$\delta$ circle, and all its closing maps commute with that
circle. Choose a circle-invariant Ehresmann connection along
one positive $w$ loop. Concatenating it six times gives the
sixth power of the same equivariant return diffeomorphism.
Changing to the invariant transport used in C.3.10 changes
this return by an equivariant isotopy: invariant horizontal
connections form an affine space, and their interpolated
flows give that isotopy.

The closing cocycle itself introduces no hidden discrepancy.
In the notation of C.3.5,
\[
 b_{h,c}=-\frac{h}{12}\Pi_c\widehat\gamma
             +\left(h\ell-\frac{h^2}{24}\right)e,
 \qquad \Pi_{c+h}\widehat\gamma
                   =\Pi_c\widehat\gamma+h e.
\]
Direct substitution gives
\[
 b_{h,c}+b_{k,c+h}=b_{h+k,c}.
\]
Thus six $h=2$ closings are the $h=12$ closing used for the
annular reduction. Returning from these coherent closings to
any stipulated R2 choices changes them by a scalar in the
connected group $\C^*$, which centralizes $S^1_\delta$;
this is again an equivariant isotopy. Every subsequent unit isotopy, finite-chart
transport and circle-gauge comparison in C.3.6 and C.3.10 is
equivariant. Consequently the annular lift and $\Phi^6$ have
the same integer $\mathfrak q$, not merely the same underlying
ordinary mapping class.

Additivity from C.3.23 now gives
\[
 6\mathfrak q(\Phi)=\mathfrak q(\Phi^6)
                   =\mathfrak q(\widehat\psi)=-J_\psi=0.
\]
Since this equality is in $\Z$, it implies
$\mathfrak q(\Phi)=0$. Its signed reduction modulo $28$
is the ordinary class $[\Phi]$, so that class is zero.
This removes the possible order-two ambiguity by an integer
argument; it does not divide the equation $6[\Phi]=0$ inside
$\Z/28$.

A smooth isotopy from the return map to the identity gives
the usual oriented product identification of its mapping
torus. A finite cyclic base change replaces the return map
by a positive power, whose class is again zero. The final
twist-data statement follows from the local family isomorphisms
and connected automorphism transitions already proved in C.4.1.
\end{proof}
\begin{proposition}[C.3.25 (Proved): an explicit eight-dimensional coboundary]
Let $\varphi$ be an orientation-preserving equivariant return map,
normalized to be the identity on a fixed tubular neighbourhood.
Let $V_\varphi$ be its actual framed-surgery six-manifold, constructed
as in C.3.14. Choose the sign of $x$ so that the unit circle
bundle of its complex line, with the boundary fibre orientation,
is the actual oriented free circle bundle. Orient its meridian
accordingly, so that $\langle x,S^2\rangle=1$. This fixes the
possible associated-line versus dual-line convention. Write
\[
 p=\langle p_1(V_\varphi)x,[V_\varphi]\rangle,
 \qquad k=\langle x^3,[V_\varphi]\rangle.
\]
There is a specified smooth spin$^c$ coboundary $Z_\varphi$ of
the actual clutching sphere $\Sigma_\varphi$. Its fourth
intersection form and characteristic covectors are
\[
 Q=\begin{pmatrix}k&1\\1&0\end{pmatrix},\qquad
 p_1(Z_\varphi)=(p+k,3),\qquad c^2=(k,1),
\]
where $c$ is the determinant class and the covectors are evaluated
on the ordered integral basis specified below. In particular
\[
 \sigma(Z_\varphi)=0,\qquad
 p_1^2=6p-3k,\qquad p_1c^2=p+k,\qquad c^4=k.
\]
With the boundary-consistent orientations, its spin$^c$ index
formula gives
\[
 \mu(\Sigma_\varphi)\equiv\frac{p-4k}{672}\pmod\Z.
\]
Boundary connected sum with $S^1\times D^7$ gives a specified
smooth spin$^c$ coboundary of the actual mapping torus, with the
same displayed characteristic numbers. This construction makes
no use of Wall's knot coordinate or the Davis generator comparison.
\end{proposition}
\begin{proof}
We first identify the boundary geometrically. Write $C$ for the
exterior of the actual fixed three-knot in the orbit six-sphere.
The framed-surgery construction is
\[
 V_\varphi=C\cup_{S^3\times S^2}(D^4\times S^2).
\]
This description, including its framing, is the direct cut-and-paste
construction of C.3.14; its later use of knot classification is not
needed here. It is spin and has $H^2=H^4=\Z$, $H^3=0$.
Let $L_x\to V_\varphi$ be the complex line with first Chern class
$x$, and put $Z_0=D(L_x)$, oriented by the base followed by its
complex line fibre. Its boundary $T=S(L_x)$ is the corresponding
principal circle bundle.

On the surgery piece $N=S^2\times D^4$, the line is the product
of the Euler-one line $\mathcal O(1)\to S^2$ and $D^4$.
Thus $T$ contains the specified product $S^3_{\rm Hopf}\times D^4$.
Attach one four-handle
\[
 D^4_{\rm Hopf}\times D^4
 \quad\hbox{along}\quad S^3_{\rm Hopf}\times D^4
\]
using this product identification and framing. Denote the resulting
eight-manifold by $Z_\varphi$. On the boundary this replaces
$S^3_{\rm Hopf}\times D^4$ by $D^4_{\rm Hopf}\times S^3$.
Over $C$, the remaining boundary is exactly the actual free circle
bundle of the clutched sphere: principal circle bundles there are
classified by their Euler class, which is the specified restriction
of $x$. On their common boundary $S^3\times S^2$, a discrepancy
between two such bundle identifications is a circle gauge with a
real logarithm, because $H^1(S^3\times S^2;\Z)=0$.
A collar cutoff corrects it to the prescribed identification.
Alexander duality gives $H^2(C;\Z)=\Z$, generated by the
meridian, so specifying this oriented Hopf restriction fixes the
entire line class on $C$. There is no additional line-bundle
parameter on the exterior. The replacement
$D^4_{\rm Hopf}\times S^3$ is now exactly the
actual fixed tubular neighbourhood with its Hopf normal action.
Hence
\[
 \partial Z_\varphi=\Sigma_\varphi
\]
as an oriented smooth manifold, with the orientation fixed by
agreement on the common free exterior. This is an identification
with the actual boundary, not merely with a sphere having a
conjectured invariant.

Explicitly, the boundary orientation of $D(L_x)$ is the positive
line-circle direction followed by the orientation of its
six-dimensional base. This is the circle-first orientation of
the actual free exterior by the defining choice of $x$.
The surgery product framing then gives the boundary-consistent
orientation on the new fixed tube. If a prior principal-connection
convention uses the dual line, replace $x$ by $-x$ throughout;
both $p$ and $k$ reverse, as does the corresponding integer
conversion. The actual zero evaluation below is unchanged.

The added four-cycle has a particularly concrete description.
Inside the displayed product region, the attachment replaces
\[
 D(\mathcal O(1))\times D^4
 \quad\hbox{by}\quad \mathbb{CP}^2\times D^4.
\]
Indeed a tubular neighbourhood of a projective line in
$\mathbb{CP}^2$ is $D(\mathcal O(1))$, and its complement is a
four-disc, with the same standard Hopf boundary identification.
Let $b$ be this positively oriented $\mathbb{CP}^2\times\{0\}$.
Its normal four-plane bundle is trivial, with the explicit
$D^4$ product framing. In particular $b^2=0$; it is not a
four-sphere with zero Pontryagin evaluation.

Let $a$ be the image of the class $\operatorname{PD}_{V_\varphi}x$
in the zero section of $D(L_x)$. A generic section of $L_x$
represents it by a closed oriented four-submanifold. Its normal
bundle in $Z_0$ is the sum of its normal line in the zero section
and the vertical line, both isomorphic to $L_x$ on that
submanifold. Therefore
\[
 a^2=\langle x^3,[V_\varphi]\rangle=k.
\]
Moreover $a\cdot b=1$: in the surgery product this intersection
is the pairing of $x$ with the meridian $S^2$, which is one by
the Euler-one convention. The four-handle exact sequence and
$H_3(V_\varphi)=0$ give
\[
 0\longrightarrow H_4(V_\varphi;\Z)
 \longrightarrow H_4(Z_\varphi;\Z)
 \longrightarrow\Z\longrightarrow0.
\]
The class $b$ uses the handle core once and maps to its generator.
Thus $(a,b)$ is an integral basis, and its intersection matrix is
the displayed $Q$. Its determinant is $-1$, so it has one
positive and one negative eigenvalue and signature zero.

The handle does not change second cohomology. Let $c$ be the
unique extension of $\pi^*x$ from $Z_0$ to $Z_\varphi$.
On $Z_0$ the tangent bundle restricts stably to
$\pi^*(TV_\varphi\oplus L_x)$, so
\[
 w_2(TZ_0)=\pi^*x\pmod2.
\]
The same degree-two handle isomorphism with $\Z/2$ coefficients
proves $w_2(TZ_\varphi)=c\pmod2$. Thus $c$ is a determinant
class of a spin$^c$ structure. It is trivial on the boundary,
which is a homotopy sphere and has its unique spin structure.
Notice that this coboundary need not be spin: $c$ is primitive.

On $a$ the stable tangent splitting gives
\[
 \langle p_1(TZ_\varphi),a\rangle=p+k,
 \qquad \langle c^2,a\rangle=k.
\]
On $b$ the normal framing gives
$TZ_\varphi|_b=T\mathbb{CP}^2\oplus\R^4$.
The class $c|_b$ is the hyperplane class, since it evaluates as
one on the original meridian projective line. Hence
\[
 \langle p_1(TZ_\varphi),b\rangle=3,
 \qquad \langle c^2,b\rangle=1.
\]
Here $p_1(T\mathbb{CP}^2)=c_1^2-2c_2=3H^2$ follows directly
from its Euler sequence. This proves the stated covectors.

Because the boundary is a homotopy sphere, the relevant relative
and absolute fourth cohomology groups agree. Their characteristic
pairing is the inverse of the integral intersection matrix:
\[
 Q^{-1}=\begin{pmatrix}0&1\\1&-k\end{pmatrix}.
\]
Multiplying the two covectors by this inverse gives exactly
\[
 p_1^2=6p-3k,\qquad p_1c^2=p+k,\qquad c^4=k.
\]

For clarity, the conversion to the boundary invariant uses only
the spin$^c$ index formula, not a knot-classification formula.
The closed eight-dimensional index density is
\[
 [e^{c/2}\widehat A]_8
 =\frac{7p_1^2-4p_2}{5760}
       -\frac{p_1c^2}{192}+\frac{c^4}{384}.
\]
Eliminating $p_2$ by the signature formula and gluing to any spin
filling of the boundary gives C.2.1's expression
\[
 \mu(\Sigma_\varphi)\equiv
 \frac{p_1^2-4\sigma}{896}
 -\frac{p_1c^2}{192}+\frac{c^4}{384}\pmod\Z.
\]
The primary normalization is
\citet[Section 2.2, equation (2.1)]{CrowleyNordstrom2020},
with its invariant divided by $28$; the spin normalization is
\citet[Section 6, equation (11)]{EellsKuiper1962}.
Substitution now gives, by exact rational arithmetic,
\[
 \frac{6p-3k}{896}-\frac{p+k}{192}+\frac{k}{384}
       =\frac{p-4k}{672}.
\]

Finally C.1.1 identifies the mapping torus with
$(S^1\times S^6)\#\Sigma_\varphi$. Boundary connected sum
of $Z_\varphi$ with $S^1\times D^7$, using a trivial determinant
line on the latter, realizes this actual boundary. Either spin
structure on the product can be used. The added summand has no
fourth cohomology and has zero signature; the displayed relative
characteristic numbers do not change.
\end{proof}

\begin{remark}[Scope and actual six-loop evaluation]
For the actual annular six-loop cap the independent geometric
calculations C.3.18--C.3.22 give
$p=-24\epsilon_{\rm cap}$ and $k=-6\epsilon_{\rm cap}$.
Here the one common sign includes the possible simultaneous
meridian-sign change just specified.
The present coboundary therefore has
\[
 p_1^2=-126\epsilon_{\rm cap},\qquad
 p_1c^2=-30\epsilon_{\rm cap},\qquad
 c^4=-6\epsilon_{\rm cap},\qquad \sigma=0.
\]
Its spin$^c$ formula is zero directly:
\[
 \epsilon_{\rm cap}
 \left(-\frac{126}{896}+\frac{30}{192}
                         -\frac{6}{384}\right)=0.
\]
This does not use the assertion $J_\psi=0$ or the Wall--Davis
conversion. By itself it evaluates the six-loop sphere. The
independent integer additivity needed for the one-loop conclusion
is supplied separately in C.3.26. The characteristic arithmetic
is checked in \texttt{d\_push\_eight\_coboundary.py}.

This is a smooth topological spin$^c$ coboundary. No complex
structure on it is asserted, and it is not identified with the
specific analytic filling requested in Part B.
\end{remark}
\begin{proposition}[C.3.26 (Proved): direct characteristic-index additivity]
For a fixed-tube-relative orbit diffeomorphism $\varphi$ of the
linear semifree sphere, form its actual surgery cap $M_\varphi$
and meridian line $L_x$ as in the geometric construction of
C.3.14. Define
\[
 \lambda(\varphi)=\frac{1}{24}
      \langle p_1(M_\varphi)x-4x^3,[M_\varphi]\rangle.
\]
This integer is additive under composition. It is unchanged by
the equivariant isotopies and fixed-tube normalizations of the
upstairs diffeomorphism, and therefore defines an integer on its
equivariant mapping class. These conclusions follow directly
from spin index theory and characteristic bordisms; they do not
use the classification or integer of three-knots, Wall's surgery
formula, or the comparison of a knot generator with a Milnor sphere.
\end{proposition}
\begin{proof}
The geometric cap is a simply connected spin six-manifold with
$H^2=H^4=\Z$ and $H^3=0$. These facts follow from the explicit
two-copy construction $N\cup_{\widetilde\varphi}(-N)$,
$N=S^2\times D^4$, by van Kampen, Mayer--Vietoris and the
unique boundary spin structure. They do not require a knot
classification. The twisted spin Dirac formula gives
\[
 \operatorname{ind}(D\otimes L_x)
 =\langle x^3/6-p_1x/24,[M_\varphi]\rangle
 =-\lambda(\varphi),
\]
so the displayed quantity is an integer.

We first prove additivity for maps fixed on the chosen tube.
Put $Y=\partial N=DP=S^2\times S^3$ and extend each map by
the identity on the second half. The meridian line on $Y$ has
a lift of each extended map: the actual upstairs circle bundle
provides it on $P$, and the identity provides it on the other
half. Their boundary values agree. Choose the lift of a
composition to be the composition of these lifts.

An oriented pair of pants has free fundamental group on its
two incoming boundary generators. The two diffeomorphisms
therefore specify a smooth $Y$-bundle over that surface, with
outgoing boundary monodromy their composition. The chosen line
bundle lifts specify a line bundle on its total space. This
total space is a compact oriented seven-manifold whose boundary
is the outgoing mapping torus minus the two incoming mapping
tori. Its stable tangent bundle restricts to their tangent
bundles plus the outward normal line. Consequently both
characteristic numbers $x^3$ and $p_1x$ are additive on these
mapping tori, by bordism invariance.

Use the explicit seven-manifold of C.3.20: a $Y$-bundle over
an annulus with monodromy the doubled map, with two
$N\times[0,1]$ pieces attached over opposite outer arcs.
The remaining outer arcs produce the twisted and identity
doubles, and the inner circle produces the mapping torus.
The chosen line lift on the annular bundle and the pullback
lines on the attached pieces agree on their attaching faces.
Thus this is an actual bordism carrying both the meridian line
and its tangent bundle, not merely a relation between abstract
intersection matrices. Its boundary orientations compare a cap
minus the identity cap with the mapping torus, with one sign
that is independent of the map. The identity cap
$S^2\times S^4$ and identity mapping torus have zero $x^3$
and zero $p_1x$. Applying the same sign to all three boundary
terms of the pair-of-pants bordism proves that $x^3$ and $p_1x$
are individually additive for fixed-tube-relative maps. Their
displayed integral combination is therefore additive as well.
In particular the proof does not infer additivity from a
three-knot group law. The line class is unambiguous here:
$H^1(Y;\Z)=0$, so changing its specified lift by a gauge does
not add a mixed mapping-torus Chern class.

It remains to check normalization choices. An orientation-
preserving equivariant sphere diffeomorphism can be made the
identity on a fixed tube by equivariant isotopies: its fixed
$S^2$ restriction is isotopic to the identity by
\citet[Theorem A]{Smale1959E}; its complex normal derivative
is homotopic to the identity since $\pi_2(U(2))=0$; and
tubular rescaling removes the remaining higher-order germ.
These are the local geometric steps, independent of any
three-knot invariant.

Two such normalizations give equivariantly diffeomorphic
clutched sphere pairs. Extend their boundary isotopy over a
collar to obtain the diffeomorphism. Compatible unitary normal
tubes along the isotopy give compatible smooth orbit structures.
Thus their orbit pairs are the same smooth pair up to
diffeomorphism. The fixed image is a three-sphere and its
oriented normal three-plane bundle is trivial. The two inherited
normal framings differ by a map $r:S^3\to SO(3)$.

We compute directly the change caused by this framing. If
$E$ is the exterior of the embedded three-sphere, its two
framed surgeries fill $E$ by $S^2\times D^4$ with boundary
identifications differing by
\[
 (a,z)\longmapsto(r(z)a,z),\qquad(a,z)\in S^2\times S^3.
\]
Cutting off the common exterior gives a characteristic bordism
between the difference of these two surgery caps and
\[
 Z_r=(S^2\times D^4)\cup_r-(S^2\times D^4),
\]
which is the oriented sphere bundle $S(\eta_r)\to S^4$ of
the rank-three bundle clutched by $r$. The meridian lines
extend over this bordism. They agree on the common boundary,
and their possible gauge discrepancy has a real logarithm
because $H^1(S^2\times S^3;\Z)=0$.

Let $\pi:Z_r\to S^4$ and choose $x$ with positive fibre
evaluation one. The oriented vertical tangent two-plane has
Euler class $2x$, since $H^2(Z_r;\Z)=\Z$ and its Euler
evaluation on the fibre is two. Its stable isomorphism
$T_{\mathrm{vert}}Z_r\oplus\R\simeq\pi^*\eta_r$ gives
\[
 4x^2=\pi^*p_1(\eta_r).
\]
Also $p_1(TS^4)=0$, so
\[
 p_1(TZ_r)=\pi^*p_1(\eta_r)=4x^2.
\]
Therefore $\langle p_1x-4x^3,[Z_r]\rangle=0$ directly.
No generator normalization of $\pi_3(SO(3))$ is needed.
The two framed-surgery indices are equal. This proves
normalization independence, and completes the direct
integer-additivity argument.
\end{proof}
\begin{corollary}[C.3.27 (Proved): a second determination by the spin$^c$ coboundary]
The actual positive-loop smooth monodromy satisfies
$[\Phi]=0\in\Theta_7$. This determination uses the explicit
eight-dimensional coboundary of C.3.25 and the direct index
additivity of C.3.26. It does not use the three-knot integer,
Wall's formula for it, the Davis generator comparison, or the
conclusion of C.3.24.
\end{corollary}
\begin{proof}
Apply C.3.26 to equivariant representatives of the actual return
maps. For the annular six-loop representative $\psi$, the two
geometric characteristic calculations, with their common
orientation sign, give
\[
 p_\psi=-24\epsilon_{\rm cap},\qquad
 k_\psi=-6\epsilon_{\rm cap}.
\]
These are the actual Pontryagin and cubic pairings proved in
C.3.18--C.3.22. Therefore the directly defined index is
\[
 \lambda(\psi)=\frac{p_\psi-4k_\psi}{24}=0
 \quad\hbox{in }\Z.
\]

We check the relation to one positive loop before reducing
modulo $28$. Choose an invariant connection on the descended
$w$-family. Its return over six consecutive positive loops is
the sixth power of its one-loop return. A positive loop changes
$c$ by $-2$. The coherent closing translations of C.3.5 obey
\[
 b_{h,c}+b_{k,c+h}=b_{h+k,c};
\]
in particular the six successive positive-$2$ closings from
$c_0-12$ to $c_0$ equal the positive-$12$ closing. Interpolating
between invariant connections gives equivariant isotopies.
The changes of transport in C.3.10, its unit and real-scalar
isotopies, and the lift comparisons of C.3.6 are equivariant
as well. Changing back to the stipulated closing choices
adds only a connected $\C^*$ automorphism commuting with the
circle action. Consequently the actual lift of $\psi$ and
$\Phi^6$ define the same equivariant mapping class.

If necessary, a lift of the orbit identity on the fixed tube
is made the identity there by its smooth circle gauge; the
gauge has a real logarithm and can be extended with a collar
cutoff. This is an equivariant isotopy and does not change the
orbit map. C.3.26 therefore applies with exactly the cap used
in the calculation above, and its normalization independence
gives
\[
 6\lambda(\Phi)=\lambda(\Phi^6)=\lambda(\psi)=0.
\]
The group here is $\Z$, so $\lambda(\Phi)=0$.

Now construct the actual eight-dimensional spin$^c$ coboundary
$Z_\Phi$ of C.3.25 for a normalized one-loop representative.
Its explicit intersection form and characteristic covectors
give, by the spin$^c$ index formula alone,
\[
 \mu(\Sigma_\Phi)
   \equiv\frac{p_\Phi-4k_\Phi}{672}
   =\frac{\lambda(\Phi)}{28}=0\pmod\Z.
\]
The positive generator in C.1.2 has invariant $1/28$, and
$\Theta_7$ has order $28$. Thus this invariant detects its
elements and $[\Sigma_\Phi]=0$. The sphere-gluing comparison
of C.1.1 gives $[\Phi]=0$. It also gives
$M_\Phi\cong S^1\times S^6$ with either spin structure having
zero invariant.
\end{proof}

\begin{remark}[Independence and shared geometric inputs]
The first determination uses the integer of equivariant
orbit knots and its Wall--Davis comparison. The present
determination uses the explicit line-disc bundle and
four-handle coboundary, the spin$^c$ index formula, and an
independent integer-additivity argument from actual
seven-dimensional characteristic bordisms. Neither conversion
nor the integer-level passage from six loops to one in this
second argument invokes the first determination. The two
methods share the proved annular representative and the actual
geometric values $p_\psi=-24\epsilon_{\rm cap}$ and
$k_\psi=-6\epsilon_{\rm cap}$. Their independence is at the
boundary-invariant and one-loop arguments; they are not
calculations from disjoint geometric data. Independent
referees and exact arithmetic checks are validation, not
additional mathematical methods.

The smooth spin$^c$ coboundary here is different from the
specific analytic filling proposed in Part B. Its construction
does not supply that filling or its almost-complex Chern data.
\end{remark}
\begin{remark}[Established surgery and index machinery]
The Hopf-handle construction used in C.3.25 is established
surgery machinery. In the proof of Lemma~3.5, Case~2,
\citet{Jiang2014E} attaches the same
$S^3_{\rm Hopf}\times D^4$ four-handle to a line-disc bundle,
expresses the resulting piece as $\mathbb{CP}^2\times D^4$,
and uses its trivial normal bundle to compute the spin$^c$
invariant. For his $k=0$ case the resulting value is
$p/(24\cdot28)$; Lemma~3.6 also exhibits the characteristic
integer $(p-4k)/24$ in his oriented decomposition conventions.
C.3.25 supplies the direct calculation with general cubic
$k$ and verifies that its boundary is the actual clutched
sphere of this family. No new general surgery construction
or invariant formula is claimed.

The relation between semifree actions, fixed three-sphere
knots, and surgery to homotopy complex projective three-spaces
is classical. The primary introduction of
\citet[pp.~125--126]{MontgomeryYang1968E} describes these
correspondences and the Pontryagin offset divisible by $24$.
The second proof here derives its needed additivity and
framing independence from the explicit bordisms and the
vertical Euler identity; it does not import the knot
classification as a premise. This identifies both the
standard literature and the exact scope of the derivation.
\end{remark}
\paragraph{The projection formula and its missing inputs.}
The current literature supplies another precise test for a proposed
orbit-knot computation. Theorem 4.1 of
\citet{GauniyalTurchin2026Root}, published online on
25 September 2026, expresses the Haefliger integer using both
double-point linking and invariants of the projected immersion.
The latter include a rationally normalized Smale invariant.
Consequently the double circle of C.3.17 alone is insufficient.
An application here would require an actual compressed projection
and its tangent-framing comparison. We opened the published
formula and its definitions, rather than importing the linking
term in isolation.

Likewise, \citet[Theorem 2.5]{Sakai2015Root} computes the
\emph{change} under a specified crossing change from oriented
double-preimage linking. To obtain an absolute value by that
route one must also identify the comparison embedding and its
class. Neither formula identifies the present trace with a
standard reference embedding. The characteristic-number cap
retains those geometric comparisons as explicit requirements.
Searches on 15 November 2026 for the six-sphere family,
its degeneration monodromy, and these knot formulas did not
locate a numerical determination for this particular family.
This states the search outcome and claims no priority.

\subsection{C.4. Change of twists and real structures}

\begin{proposition}[C.4.1 (Proved): twist independence in the sphere cases]
For the fixed multiplicities $(3,4)$ and all admissible data in R5
with $|p|=1$, the monodromy class of the positive $w$-circle is the
same class in $\Theta_7$, with the complex orientation on each fibre.
\end{proposition}

\begin{proof}
The explicit sufficiency maps in E.1 depend holomorphically on a
small $c$-disc: the lattice defects and toric integral exponents are
constant, the periods and unit factors are holomorphic, and the
translation correction is solved in the fixed sheaf $V$ with
$H^1(B,V)=0$. This is the parameter-dependent assertion of E.1.3,
not just the fibrewise classification. For the second twist
coefficient $s_2$, the parameter identification is
$W=e^{-\pi i(c+s_2)}=(-1)^{s_2}w$.
The rotation of the punctured disc preserves its positive
fundamental-group generator. Local family identifications differ
by elements of $\Aut(X)=\C^*$, a connected group, so their
transition changes are smoothly isotopic to the identity. Following
these identifications around the circle conjugates the monodromy
in the oriented mapping-class group. Since $\Theta_7$ is abelian,
the resulting class is unchanged. If the lattice sign in R5 is
$-1$, the comparison map is nevertheless holomorphic and preserves
the complex orientation; no extra sign is introduced.
The common class is zero by C.3.24. This last numerical consequence
uses only the already proved equality of the classes, not an
additional deformation or conjugation argument.
\end{proof}

\begin{proposition}[C.4.2 (Proved): conjugation gives no torsion constraint]
The antiholomorphic identification in R6 covers
$w\mapsto-\bar w$. At a fixed base point on the imaginary
$w$-axis it identifies the positive-loop monodromy with the
inverse-loop monodromy by an orientation-reversing fibre map.
Its action on $\Theta_7$ is negation, so the resulting equation is
the tautology $-[\Phi]=-[\Phi]$, not $2[\Phi]=0$.
\end{proposition}

\begin{proof}
From $c'=1-\bar c$ we obtain
$e^{-\pi ic'}=-\overline{e^{-\pi ic}}$.
Complex conjugation reverses the parameter circle. An
antiholomorphic map in complex dimension three reverses the real
orientation because its determinant has sign $(-1)^3$.
Conjugating a sphere clutching map by a reflection changes the
orientation of its glued homotopy sphere, hence negates its
$\Theta_7$ class. Equivalently the Eells--Kuiper invariant changes
sign. Applying this to the inverse loop proves the assertion.
This confirms the last sentence of Lead 7 without extracting a
spurious order-two restriction.
\end{proof}

\begin{corollary}[C.5.1 (Refuted): the specified analytic filling premise]
The analytic descended filling and its resolution $W_8$
requested in Parts B.4 and C.2 do not exist. Consequently
there are no values of $c_1,c_2,p_1,p_1^2,p_1c_1^2,c_1^4$,
relative lifts, intersection form or signature for that
specified analytic object. The actual positive-loop
monodromy is nevertheless
\[
                    [\Phi]=0\in\Theta_7\simeq\Z/28.
\]
\end{corollary}
\begin{proof}
A descended analytic filling of the specified family
would be a proper analytic filling before any resolution,
and its finite pullback would remain proper. B.4.33
excludes both possibilities. Characteristic numbers of
an unrelated manifold cannot supply the requested values.

C.3.24 determines the actual monodromy by the equivariant
orbit-knot integer. C.3.25--C.3.27 construct an actual
smooth spin$^c$ coboundary of the clutching sphere and
mapping torus and give a second determination by the
index formula. Both use the verified actual pairings
$p_\psi=-24\epsilon_{\rm cap}$ and
$k_\psi=-6\epsilon_{\rm cap}$ and pass from six loops
to one in an integer group before reduction modulo $28$.
Their distinct boundary-invariant arguments and shared
geometric inputs are stated in C.3.27. C.4.1 gives the
same zero for all the compared sphere twist data, and
C.4.2 gives the antiholomorphic transformation law.

The smooth coboundary of C.3.25 is fully specified and
has the intersection and spin$^c$ characteristic data
computed there. No complex structure on it is asserted,
and its determinant class is not identified with the first
Chern class of an almost-complex tangent bundle. This
distinction is essential now that the prescribed analytic
construction has been refuted.
\end{proof}
